% Copied on 2026-11-03 from eth-ainit-hs26/w1-lecture-material, commit 9792819, file style/ainit-preamble.tex. To pull a fix, diff against the source and recopy.
% ===== ainit-preamble.tex: the official deck style of AI in Industry (AINIT), HS26 =====
% "Codex plates". Chosen by Carlos on 9 October 2026 from a two round style board. Round one
% had Codex, Atlas, Spec sheet, Sheet and Cartridge; round two had Codex violet, Codex plates,
% Vellum atlas and Cool atlas. CLAUDE.md in this folder records the decision, the rules and the
% traps. Read it before changing anything here.
%
% What it is. Every slide is a leaf of a codex: warm vellum inside a fine double ruled block,
% pricking dots in the outer margin, the title above a double rule, a violet folio and an italic
% slug in the lower margin. The figures speak a technical grammar in fine ink: sharp rectangles,
% 45 degree hatching in violet lines with no solid fill, dashed outlines for whatever is refused
% or not taken, right angled connectors, numbered balloons on hairline leaders, and a large
% numeral in a heavy square. Violet is the rubric: the folio, key words, the chosen path, the
% current weekend. Ochre is the second accent. Ge'ez numerals number the games, the balloons and
% the sections, from one to five. Arabic digits take over beyond five, because five is how many
% numerals the room learns in the Ge'ez games.
%
% Each deck does:
%   \documentclass[aspectratio=169]{beamer}
%   % Copied on 2026-11-03 from eth-ainit-hs26/w1-lecture-material, commit 9792819, file style/ainit-preamble.tex. To pull a fix, diff against the source and recopy.
% ===== ainit-preamble.tex: the official deck style of AI in Industry (AINIT), HS26 =====
% "Codex plates". Chosen by Carlos on 9 October 2026 from a two round style board. Round one
% had Codex, Atlas, Spec sheet, Sheet and Cartridge; round two had Codex violet, Codex plates,
% Vellum atlas and Cool atlas. CLAUDE.md in this folder records the decision, the rules and the
% traps. Read it before changing anything here.
%
% What it is. Every slide is a leaf of a codex: warm vellum inside a fine double ruled block,
% pricking dots in the outer margin, the title above a double rule, a violet folio and an italic
% slug in the lower margin. The figures speak a technical grammar in fine ink: sharp rectangles,
% 45 degree hatching in violet lines with no solid fill, dashed outlines for whatever is refused
% or not taken, right angled connectors, numbered balloons on hairline leaders, and a large
% numeral in a heavy square. Violet is the rubric: the folio, key words, the chosen path, the
% current weekend. Ochre is the second accent. Ge'ez numerals number the games, the balloons and
% the sections, from one to five. Arabic digits take over beyond five, because five is how many
% numerals the room learns in the Ge'ez games.
%
% Each deck does:
%   \documentclass[aspectratio=169]{beamer}
%   % Copied on 2026-11-03 from eth-ainit-hs26/w1-lecture-material, commit 9792819, file style/ainit-preamble.tex. To pull a fix, diff against the source and recopy.
% ===== ainit-preamble.tex: the official deck style of AI in Industry (AINIT), HS26 =====
% "Codex plates". Chosen by Carlos on 9 October 2026 from a two round style board. Round one
% had Codex, Atlas, Spec sheet, Sheet and Cartridge; round two had Codex violet, Codex plates,
% Vellum atlas and Cool atlas. CLAUDE.md in this folder records the decision, the rules and the
% traps. Read it before changing anything here.
%
% What it is. Every slide is a leaf of a codex: warm vellum inside a fine double ruled block,
% pricking dots in the outer margin, the title above a double rule, a violet folio and an italic
% slug in the lower margin. The figures speak a technical grammar in fine ink: sharp rectangles,
% 45 degree hatching in violet lines with no solid fill, dashed outlines for whatever is refused
% or not taken, right angled connectors, numbered balloons on hairline leaders, and a large
% numeral in a heavy square. Violet is the rubric: the folio, key words, the chosen path, the
% current weekend. Ochre is the second accent. Ge'ez numerals number the games, the balloons and
% the sections, from one to five. Arabic digits take over beyond five, because five is how many
% numerals the room learns in the Ge'ez games.
%
% Each deck does:
%   \documentclass[aspectratio=169]{beamer}
%   % Copied on 2026-11-03 from eth-ainit-hs26/w1-lecture-material, commit 9792819, file style/ainit-preamble.tex. To pull a fix, diff against the source and recopy.
% ===== ainit-preamble.tex: the official deck style of AI in Industry (AINIT), HS26 =====
% "Codex plates". Chosen by Carlos on 9 October 2026 from a two round style board. Round one
% had Codex, Atlas, Spec sheet, Sheet and Cartridge; round two had Codex violet, Codex plates,
% Vellum atlas and Cool atlas. CLAUDE.md in this folder records the decision, the rules and the
% traps. Read it before changing anything here.
%
% What it is. Every slide is a leaf of a codex: warm vellum inside a fine double ruled block,
% pricking dots in the outer margin, the title above a double rule, a violet folio and an italic
% slug in the lower margin. The figures speak a technical grammar in fine ink: sharp rectangles,
% 45 degree hatching in violet lines with no solid fill, dashed outlines for whatever is refused
% or not taken, right angled connectors, numbered balloons on hairline leaders, and a large
% numeral in a heavy square. Violet is the rubric: the folio, key words, the chosen path, the
% current weekend. Ochre is the second accent. Ge'ez numerals number the games, the balloons and
% the sections, from one to five. Arabic digits take over beyond five, because five is how many
% numerals the room learns in the Ge'ez games.
%
% Each deck does:
%   \documentclass[aspectratio=169]{beamer}
%   \input{../style/ainit-preamble}      % this file
%   \input{../style/ainit-course}        % the facts of the course and the three weekends frame
%   \renewcommand{\deckslug}{AI in Industry, weekend 1}
% and is compiled with lualatex FROM the deck's own directory, three times: the leaf, the dots
% and the balloons are placed with remember picture and settle on the third pass. Fonts are
% EB Garamond, Garamond Math and Libertinus Mono (all in TeX Live) and Kefa III for Ethiopic
% (a macOS system font; set \geezfontpath before this file to use another folder).
%
% What a deck sets:
%   \title \subtitle \author \institute \date       the title page
%   \renewcommand{\deckslug}{...}                   the italic slug in the footer
%   \def\titlegeez{...}                             optional: a Ge'ez line above the cover title
%   \plateImage \plateCredit \quoteText \quoteCitation   the cold open, shown by \coldopenframe
% Frames 'title' and 'cold open' are [plain,noframenumbering]; a \section makes its own divider.
%
% The kit a deck may use:
%   macros   \hl{key words}  \geez{...}  \geezNumeral{n}  \cxballoon  \treelink  \weblink  \qrlink
%            \framedimgh{height}{file}  \coldopenframe  \begin{codeplate}...\end{codeplate}
%            \begin{cxlist}[\small]...\end{cxlist}  a denser list for reference slides
%   styles   chip, chip strong, chip dim, edge strong, edge dim, elab, elab dim, axis, vec a,
%            vec b, vlab a, vlab b, anglearc, anglelab, gloss, hatch, hatch fine
%   colours  paper ink muted brand accent2 good bad ruling ruledark nib
%
% House rules that bind every deck (the lecture-slides skill has the long form):
%   light and projection first; no red against green (good is blue, bad is vermillion); no ALL
%   CAPS; no dash or double hyphen in prose; nothing under a diagram; a section divider carries a
%   numeral, "Section N", the title and a band, nothing else; only the title page names the
%   lecturer; the only number in the chrome is the folio; no cell grids, column letters, tabs,
%   title blocks or arcade chrome. Never name a colour fg or bg: beamer rebinds both.

\usepackage{fontspec}
\usepackage{amsmath}
\usepackage{booktabs}
\usepackage{unicode-math}
\usefonttheme{professionalfonts}
\usefonttheme{serif}
\usepackage{tikz}
\usetikzlibrary{arrows.meta, positioning, calc, backgrounds, fit, overlay-beamer-styles,
                decorations.pathreplacing, patterns.meta}
\setbeamertemplate{navigation symbols}{}

% === type: EB Garamond (TeX Live), Garamond-Math for the maths, Kefa III for Ethiopic ===
\setmainfont{EBGaramond}[
  Extension      = .otf,
  UprightFont    = *-Medium,
  ItalicFont     = *-MediumItalic,
  BoldFont       = *-SemiBold,
  BoldItalicFont = *-SemiBoldItalic,
  Ligatures      = TeX,
  Numbers        = Lining,
  Scale          = 1.08]
\setmathfont{Garamond-Math.otf}[Scale=1.08]
% a monospaced face for code and addresses, sized to the lowercase of the main face
\setmonofont{LibertinusMono-Regular.otf}[Scale=MatchLowercase]

% Kefa III is a macOS system font. Its light and bold instances are variable font instances.
% Without the file, the deck still compiles: Ge'ez text then lacks glyphs and a warning says so.
\providecommand{\geezfontpath}{/System/Library/Fonts/Supplemental/}
\IfFileExists{\geezfontpath KefaIII.ttf}{%
  \newfontfamily\geezfont[Path=\geezfontpath]{KefaIII.ttf}
  % the light instance, for large numerals
  \newfontfamily\geezlightfont[Path=\geezfontpath, RawFeature={instance=Light}]{KefaIII.ttf}
  % the bold instance, for the numerals in balloons (the regular weight is too thin there)
  \newfontfamily\geezboldfont[Path=\geezfontpath, RawFeature={instance=Bold}]{KefaIII.ttf}
}{%
  \PackageWarningNoLine{ainit-preamble}{Kefa III was not found in \geezfontpath. Ge'ez text
    will lack glyphs. Set \string\geezfontpath\space before this file.}%
  \newcommand{\geezfont}{}\newcommand{\geezlightfont}{}\newcommand{\geezboldfont}{}}
\newcommand{\geez}[1]{{\geezfont #1}}
\newcommand{\geezlight}[1]{{\geezlightfont #1}}

% === colours. Never name one fg or bg (beamer rebinds both). ===
\definecolor{paper}{HTML}{F9F4E3}     % vellum, lightness above 93 percent
\definecolor{ink}{HTML}{241F1A}       % warm black, iron gall
\definecolor{muted}{HTML}{5E5448}     % secondary text, darker than black!55
\definecolor{brand}{HTML}{58409A}     % violet in the rubric role
\definecolor{accent2}{HTML}{8F5A06}   % ochre, the second accent
\definecolor{good}{HTML}{0072B2}      % Okabe-Ito blue: a win
\definecolor{bad}{HTML}{D55E00}       % Okabe-Ito vermillion: a loss. No red against green, ever.
\colorlet{ruling}{ink!40!paper}       % the ruled lines of the leaf
\colorlet{ruledark}{ink!62!paper}     % the heavier ruled lines
\colorlet{nib}{ink!88!paper}          % the outline ink of the figures

\setbeamercolor{normal text}{fg=ink, bg=paper}
\setbeamercolor{structure}{fg=brand}
\setbeamercolor{alerted text}{fg=brand}
\setbeamersize{text margin left=15mm, text margin right=15mm}

% === links: violet, and a raw address keeps the mono face, so a link is marked twice ===
% Beamer loads hyperref. The border is cleared; colorlinks stays off because it would also tint
% the section title on the divider.
\hypersetup{pdfborder={0 0 0}, breaklinks=true}
\urlstyle{tt}
\renewcommand{\UrlFont}{\ttfamily\color{brand}}
% \weblink{https://...}{the words the room reads}: an inline link whose visible text is prose.
\newcommand{\weblink}[2]{\href{#1}{\color{brand}#2}}

% \qrlink{2.6cm}{https://...}: a vector QR on a plain white plate in a hairline mount. The code is
% generated from the address at compile time, so it cannot disagree with the address printed
% beside it. Always print the address next to it.
% Two guards, both invisible in the source. nolink: qrcode makes the whole code a link by default,
% and a link is painted in the link colour. \color{black}: qrcode draws in the current colour, so
% on a full bleed frame where \color{white} is active the modules would vanish.
\usepackage{qrcode}
\newcommand{\qrlink}[2]{%
  \tikz{\node[inner sep=4pt, fill=white, draw=ruledark, line width=0.6pt, line join=miter]
    {\color{black}\qrcode[nolink, height=#1]{#2}};}}

% \framedimgh{2.5cm}{img/photo.jpg}: a photograph or figure of the given height in a hairline mount,
% standing 1.4 mm off the picture as the plate of the cold open does.
\newcommand{\framedimgh}[2]{%
  \tikz{\node[inner sep=1.4mm, draw=ruling, line width=0.4pt, line join=miter]
    {\includegraphics[height=#1]{#2}};}}

% === the overlay style used by every piece of page furniture ===
\tikzset{cx page/.style={remember picture, overlay, shift={(current page.north west)},
  x=1mm, y=-1mm}}

% === hatching: the one fill of the figures. Violet lines at 45 degrees, no solid colour under them ===
\tikzset{
  hatch/.style={pattern={Lines[angle=45, distance=0.95mm, line width=0.42pt]}, pattern color=brand},
  hatch fine/.style={pattern={Lines[angle=45, distance=0.62mm, line width=0.36pt]}, pattern color=brand},
}

% === Ge'ez numerals up to five, digits beyond ===
% The room learns ፩ to ፭ in the Ge'ez games (the course numbers the games with them).
\newcommand{\geezNumeral}[1]{\ifcase#1 0\or ፩\or ፪\or ፫\or ፬\or ፭\else\number#1\fi}
\newcommand{\cxnum}[1]{\ifnum#1<6\relax{\geezboldfont\geezNumeral{#1}}\else\number#1\fi}
\newcounter{glossn}

% a box that is scaled down to fit when its natural width exceeds #1, and left alone otherwise
\newsavebox{\cxfitbox}
\newcommand{\cxfit}[2]{%
  \sbox{\cxfitbox}{#2}%
  \ifdim\wd\cxfitbox>#1\relax\resizebox{#1}{!}{\usebox{\cxfitbox}}\else\usebox{\cxfitbox}\fi}

% === the footer slug: each deck renews it ===
\providecommand{\deckslug}{AI in Industry}

% === the leaf: ruled block with a fine double rule, pricking marks in the outer margin ===
\setbeamertemplate{background}{%
  \setcounter{glossn}{0}%
  \begin{tikzpicture}[cx page]
    \draw[ruling, line width=0.55pt] (8,5) rectangle (152,82);
    \draw[ruling, line width=0.25pt] (8.9,5.9) rectangle (151.1,81.1);
    \foreach \i in {0,...,12} {\fill[ruling] (3.2,{5+6.4167*\i}) circle (0.3);}
  \end{tikzpicture}}

% === the title sits above the first ruled line of the block; an over long title is scaled to fit ===
\setbeamertemplate{frametitle}{%
  \vbox to 17.5mm{\offinterlineskip
    \vskip 8.3mm
    \hbox to\textwidth{\vrule height 5.3mm depth 1.9mm width 0pt
      \cxfit{\textwidth}{\color{ink}\fontsize{19}{22}\selectfont\bfseries\insertframetitle}\hfil}%
    \vskip 0.8mm
    \hbox to\textwidth{\color{ruledark}\vrule height 0.7pt depth 0pt width\textwidth}%
    \vskip 0.6mm
    \hbox to\textwidth{\color{ruling}\vrule height 0.3pt depth 0pt width\textwidth}%
    \vss}}

% === the folio in the lower margin, the italic slug beside it ===
\setbeamertemplate{footline}{%
  \vbox to 8.2mm{\offinterlineskip\vss
    \hbox to\paperwidth{\hskip 15mm
      {\color{brand}\fontsize{12}{14}\selectfont\insertframenumber}\hskip 3.4mm
      {\color{muted}\footnotesize\itshape\deckslug}\hfil}%
    \vskip 2.2mm}}

% === the title page ===
% \titlegeez (a Ge'ez line over the title) and \geezYear (the year, in Ge'ez numerals, after the
% date) are optional. \geezYear comes from ainit-course.tex.
\setbeamertemplate{title page}{%
  \begin{tikzpicture}[cx page]
    \node[text=brand, font=\fontsize{17}{20}\selectfont, anchor=base] at (80,20.5)
      {\ifdefined\titlegeez\geez{\titlegeez}\fi};
    \node[text=ink, font=\fontsize{27}{30}\selectfont\bfseries, text width=104mm, align=center,
      anchor=north, inner sep=0] at (80,25) {\inserttitle};
    \draw[ruledark, line width=0.8pt] (62,50.5) -- (98,50.5);
    \draw[ruling, line width=0.3pt] (62,51.5) -- (98,51.5);
    \node[text=muted, font=\large\itshape, anchor=north, inner sep=0] at (80,55)
      {\insertsubtitle};
    \node[text=ink, font=\large, anchor=base, inner sep=0] at (80,67) {\insertauthor};
    \node[text=muted, font=\normalsize\itshape, anchor=base, inner sep=0] at (80,72.2)
      {\insertinstitute};
    \node[text=muted, font=\normalsize\itshape, anchor=base, inner sep=0] at (80,77.2)
      {\insertdate\ifdefined\geezYear\hspace{3.2mm}{\upshape\color{brand}\geez{\geezYear}}\fi};
  \end{tikzpicture}}

% === the section page: a numeral in a heavy square, Section N, the title, a hatched band ===
% The numeral is Ge'ez up to five and a digit beyond. A title that is wider than 92 mm at the
% display size wraps over several lines at a smaller size, and the band moves down under it.
\newif\ifcxwrap
\setbeamertemplate{section page}{%
  \sbox{\cxfitbox}{\fontsize{36}{40}\selectfont\bfseries\insertsectionhead}%
  \ifdim\wd\cxfitbox>92mm\relax\cxwraptrue\else\cxwrapfalse\fi
  \begin{tikzpicture}[cx page]
    \draw[brand, line width=1.8pt, line join=miter] (17,30.2) rectangle (44,57.2);
    \draw[brand, line width=0.4pt] (18.9,32.1) rectangle (42.1,55.3);
    \ifnum\insertsectionnumber<6\relax
      \node[text=brand, font=\fontsize{54}{54}\selectfont, anchor=center] at (30.56,43.7)
        {\geezlight{\geezNumeral{\insertsectionnumber}}};
    \else
      \node[text=brand, font=\fontsize{54}{54}\selectfont, anchor=center] at (30.56,43.7)
        {\number\insertsectionnumber};
    \fi
    \node[text=muted, font=\fontsize{14}{16}\selectfont\itshape, anchor=base west, inner sep=0]
      at (53,33.2) {Section \insertsectionnumber};
    \ifcxwrap
      \node[text=ink, font=\fontsize{26}{29}\selectfont\bfseries, anchor=north west, inner sep=0,
        text width=92mm, align=flush left] (cxst) at (53,36.5) {\insertsectionhead};
      \fill[hatch] ([shift={(0,3.4)}]cxst.south west) rectangle ([shift={(92,7)}]cxst.south west);
      \draw[brand, line width=0.5pt] ([shift={(0,3.4)}]cxst.south west) -- ++(92,0)
        ([shift={(0,7)}]cxst.south west) -- ++(92,0);
    \else
      \node[text=ink, font=\fontsize{36}{40}\selectfont\bfseries, anchor=base west, inner sep=0]
        at (53,47.5) {\insertsectionhead};
      \fill[hatch] (53,53.6) rectangle (145,57.2);
      \draw[brand, line width=0.5pt] (53,53.6) -- (145,53.6) (53,57.2) -- (145,57.2);
    \fi
  \end{tikzpicture}}
\AtBeginSection[]{\begin{frame}[plain,noframenumbering]\usebeamertemplate{section page}\end{frame}}

% === rubrication: key words in semibold violet ===
\newcommand{\hl}[1]{\textcolor{brand}{\bfseries #1}}

% === the cold open: the plate in a hairline mount, the verse, the citation ===
% The deck sets \plateImage (a file name), \plateCredit, \quoteText (lines separated by \\, key
% words in \hl{}) and \quoteCitation, and calls \coldopenframe inside a [plain] frame.
\newcommand{\coldopenframe}{%
  \begin{tikzpicture}[cx page]
    \node[anchor=north west, inner sep=0] (plate) at (16,15)
      {\includegraphics[height=52mm]{\plateImage}};
    \draw[ruling, line width=0.4pt] ([shift={(-1.6,-1.6)}]plate.north west) rectangle
      ([shift={(1.6,1.6)}]plate.south east);
    \node[anchor=north, inner sep=0, text=muted, font=\scriptsize\itshape]
      at ($(plate.south)+(0,3.2)$) {\plateCredit};
    \node[anchor=west, inner sep=0, text width=86mm, align=flush left,
      font=\fontsize{11.5}{17}\selectfont] at (62,43.5) {%
      \quoteText\par\vskip 4mm
      {\color{ruling}\vrule height 0.4pt depth 0pt width 22mm}\par\vskip 2mm
      {\small\itshape\color{muted}\quoteCitation}};
  \end{tikzpicture}}

% === a numbered balloon: a hairline circle with a numeral, centre (#1,#2) in the picture's units ===
\newcommand{\cxballoon}[4]{% x, y, radius in mm, numeral
  \draw[brand, line width=0.7pt, fill=paper] (#1,#2) circle [radius=#3mm];
  \node[text=brand, inner sep=0, font=\fontsize{14}{14}\selectfont] at (#1,#2) {\cxnum{#4}};}

% \treelink{edge style}{label style}{from}{to}{label}: one right angled branch of a tree.
% It leaves the east anchor of the parent, runs level to a junction at two fifths of the gap,
% turns once to the height of the child and arrives level at its west anchor. The first branch
% drawn from a parent carries the shared stretch, so declare the chosen branch first: later
% branches from the same parent begin at the junction, and a dashed branch never overprints the
% heavier path.
\newcommand{\treelink}[5]{%
  \ifcsname cx@trunk@#3\endcsname
    \draw[#1] let \p1=(#3.east), \p2=(#4.west), \p3=($(\p1)!0.4!(\p2)$),
                  \p4=($(\x3,\y2)!0.5!(\p2)$) in
      (\p4) node[#2, above=0.4mm] {#5}
      (\x3,\y1) -- (\x3,\y2) -- (\p2);
  \else
    \expandafter\def\csname cx@trunk@#3\endcsname{}%
    \draw[#1] let \p1=(#3.east), \p2=(#4.west), \p3=($(\p1)!0.4!(\p2)$),
                  \p4=($(\x3,\y2)!0.5!(\p2)$) in
      (\p4) node[#2, above=0.4mm] {#5}
      (\p1) -- (\x3,\y1) -- (\x3,\y2) -- (\p2);
  \fi}

% === list markers: a small square hatched in violet ===
\newcommand{\cxhatchsq}[1]{%
  \raisebox{0.12ex}{\tikz[baseline=0pt]{%
    \fill[hatch fine] (0,0) rectangle (#1,#1);
    \draw[brand, line width=0.55pt, line join=miter] (0,0) rectangle (#1,#1);}}}
\setbeamertemplate{itemize item}{\cxhatchsq{3mm}}
\setbeamertemplate{itemize subitem}{\cxhatchsq{2.1mm}}
\setbeamertemplate{enumerate item}{%
  \tikz[baseline=(enumn.base)]\node[draw=brand, line width=0.7pt, circle, inner sep=0.5pt,
    minimum size=5.4mm, font=\small, text=brand, fill=paper] (enumn) {\cxnum{\insertenumlabel}};}
\setbeamertemplate{enumerate subitem}{\insertenumlabel.}
\setbeamerfont{itemize/enumerate body}{size=\Large}
\setbeamerfont{itemize/enumerate subbody}{size=\large}
\setbeamertemplate{itemize/enumerate body begin}{\vspace*{0mm}}
% beamer resets the list parameters inside \list, so the spacing is set in the list defaults
\makeatletter
\def\@listi{\leftmargin\leftmargini
            \topsep 0\p@ \parsep 0\p@
            \itemsep 6mm \@plus1mm \@minus1mm}
\let\@listI\@listi
\def\@listii{\leftmargin\leftmarginii
             \topsep 0\p@ \parsep 0\p@
             \itemsep 2mm \@plus0.5mm \@minus0.5mm}
\makeatother
% a denser list for reference slides (administration, rosters, small print): a smaller marker and
% the text size #1, \small by default. Set \itemsep right after \begin{cxlist} as with itemize.
\newenvironment{cxlist}[1][\small]{%
  \begingroup
  \setbeamertemplate{itemize item}{\cxhatchsq{2.2mm}}%
  \setbeamerfont{itemize/enumerate body}{size=#1}%
  \begin{itemize}}{\end{itemize}\endgroup}

% === the code plate: mono text in a hairline box under a hatched band ===
% Use it in a [fragile] frame. It is built for lines up to about 60 characters: a longer line
% runs past the box, so shorten the code or the prompt instead of shrinking the type.
\usepackage{listings}
\usepackage[skins,listings]{tcolorbox}
\lstdefinestyle{ainit}{
  basicstyle=\ttfamily\normalsize, keywordstyle=\color{brand}\bfseries,
  commentstyle=\color{muted}\itshape, stringstyle=\color{accent2},
  showstringspaces=false, columns=fullflexible, keepspaces=true, breaklines=false,
  aboveskip=0pt, belowskip=0pt, language=Python}
\newtcblisting{codeplate}[1][]{%
  enhanced, listing only, listing options={style=ainit},
  colback=paper, colframe=ruledark, boxrule=0.6pt, sharp corners,
  left=4mm, right=4mm, top=5.6mm, bottom=2.2mm, before skip=1mm, after skip=1mm,
  overlay unbroken={%
    \fill[hatch] ([yshift=-2.9mm]frame.north west) rectangle (frame.north east);
    \draw[brand, line width=0.5pt]
      ([yshift=-2.9mm]frame.north west) -- ([yshift=-2.9mm]frame.north east);},
  #1}

% === the TikZ styles of the figures ===
\tikzset{
  >={Triangle[length=2.5mm, width=1.9mm]},
  chip/.style={draw=nib, line width=0.7pt, line join=miter, minimum width=1.85cm,
    minimum height=0.82cm, inner sep=3pt, text=ink, fill=paper, font=\large},
  chip strong/.style={chip, draw=brand, line width=1.0pt,
    path picture={%
      \fill[hatch, even odd rule]
        (path picture bounding box.south west) rectangle (path picture bounding box.north east)
        ([shift={(1.25mm,1.25mm)}]path picture bounding box.south west) rectangle
        ([shift={(-1.25mm,-1.25mm)}]path picture bounding box.north east);
      \draw[brand, line width=0.35pt]
        ([shift={(1.25mm,1.25mm)}]path picture bounding box.south west) rectangle
        ([shift={(-1.25mm,-1.25mm)}]path picture bounding box.north east);}},
  chip dim/.style={chip, draw=muted, line width=0.6pt, dash pattern=on 2.4pt off 1.8pt,
    text=muted},
  edge strong/.style={draw=brand, line width=1.3pt, line join=miter,
    -{Triangle[length=2.9mm, width=2.2mm]}},
  edge dim/.style={draw=muted, line width=0.6pt, line join=miter, dash pattern=on 3pt off 2.2pt,
    -{Triangle[length=2.5mm, width=1.8mm]}},
  elab/.style={font=\small\itshape, text=brand, fill=paper, inner sep=1.5pt},
  elab dim/.style={font=\small\itshape, text=muted, fill=paper, inner sep=1.5pt},
  axis/.style={draw=muted, line width=0.6pt, -{Triangle[length=2.3mm, width=1.6mm]}},
  vec a/.style={draw=brand, line width=1.4pt, -{Triangle[length=3.1mm, width=2.3mm]}},
  vec b/.style={draw=accent2, line width=1.4pt, -{Triangle[length=3.1mm, width=2.3mm]}},
  vlab a/.style={text=brand, font=\normalsize},
  vlab b/.style={text=accent2, font=\normalsize},
  anglearc/.style={draw=ink, line width=0.6pt,
    {Triangle[length=2mm, width=1.4mm]}-{Triangle[length=2mm, width=1.4mm]}},
  anglelab/.style={text=ink, font=\normalsize},
  % a note: a numbered balloon on a hairline leader, with its text. The balloons number
  % themselves 1, 2, 3 on each page.
  gloss/.style={font=\small\itshape, text=muted, text width=3.9cm, align=flush left,
    inner xsep=6.9mm, inner ysep=0.6mm,
    execute at begin node={\stepcounter{glossn}},
    path picture={%
      \draw[brand, line width=0.5pt]
        ($(path picture bounding box.north west)+(2.9mm,-5.4mm)$) --
        ($(path picture bounding box.south west)+(2.9mm,0.3mm)$);
      \draw[brand, line width=0.7pt, fill=paper]
        ($(path picture bounding box.north west)+(2.9mm,-2.9mm)$) circle (2.5mm);
      \pgftext[at={\pgfpointadd{\pgfpointanchor{path picture bounding box}{north west}}%
        {\pgfpoint{2.9mm}{-2.9mm}}}]
        {\color{brand}\fontsize{10.5}{10.5}\selectfont\upshape\cxnum{\value{glossn}}}}},
}
      % this file
%   % Copied on 2026-11-03 from eth-ainit-hs26/w1-lecture-material, commit 9792819, file style/ainit-course.tex. To pull a fix, diff against the source and recopy.
% ===== ainit-course.tex: the facts of AINIT HS26 and the three weekends frame =====
% Load this after ainit-preamble.tex in any deck that shows the course roadmap or the title page's
% Ge'ez year. The facts live here once, so a change of date or title is one edit.
%
% Where the facts come from. Dates, themes and titles are the master spreadsheet's Syllabus tab and
% agree with the course website. The three project names differ between the website, the sheet and
% the introduction deck. This file follows the introduction deck (admin-deck/), because that is the
% deck the participants see first. The topics line of each weekend is a short reading of the
% weekend's slots. Re-read the sheet before the semester starts and after any change of schedule.
%
% For the next academic year change \geezYear, the dates, and whatever else moved.

% 2026 in Ge'ez numerals: 20 x 100 + 20 + 6. The title page prints it after the date.
\newcommand{\geezYear}{፳፻፳፮}

% ---- the three weekends (a, b, c): number, title, dates, topics, project, Faust epithet ----
% The epithets come from the course's cold open: Goethe's Faust renders the Greek word logos as
% the Word, the Thought and the Act, and refuses the Power. They map to the three weekends.
\newcommand{\wkaNum}{1}\newcommand{\wkaTitle}{Large language models}
\newcommand{\wkaDates}{16 and 17 October}
\newcommand{\wkaTopics}{Ge'ez games, prompting, OPRO, MLOps}
\newcommand{\wkaProject}{LLM-based marketing agent}\newcommand{\wkaMuse}{the Word}
\newcommand{\wkbNum}{2}\newcommand{\wkbTitle}{Representation learning}
\newcommand{\wkbDates}{13 and 14 November}
\newcommand{\wkbTopics}{Recommendation systems, matrix factorization, bandits}
\newcommand{\wkbProject}{Recommendation systems}\newcommand{\wkbMuse}{the Thought}
\newcommand{\wkcNum}{3}\newcommand{\wkcTitle}{Reinforcement learning}
\newcommand{\wkcDates}{11 and 12 December}
\newcommand{\wkcTopics}{MDPs, Bellman equations, SARSA, REINFORCE}
\newcommand{\wkcProject}{Reinforcement learning for delivery optimization}
\newcommand{\wkcMuse}{the Act}

% ===== frame: the three weekends, three plates side by side =====
% The current weekend has a hatched header band and an ochre flag. Set it with
% \renewcommand{\currentweekend}{b} in a weekend 2 deck. The refused Power hangs as a struck label
% above the gap between the second and the third plate, with a dashed phantom line running down that
% gap. Switch it off with \cxpowerfalse after loading this file.
\newif\ifcxpower
\cxpowertrue
\newcommand{\cxW}{4.08}   % plate width in cm
\newcommand{\cxG}{0.38}   % gap between plates in cm
\newcommand{\cxH}{5.0}    % plate height in cm
\newcommand{\cxtop}{0.56} % height of the row of flags above the plates, in cm
\newcommand{\cxband}{0.44}% height of the header band of a plate, in cm

% the capital form of a Muse: "the Word" becomes "The Word"
\def\cxcaptheA the #1\relax{The #1}
\newcommand{\cxMuse}[1]{{\edef\cxtmp{\csname wk#1Muse\endcsname}\expandafter\cxcaptheA\cxtmp\relax}}

% one plate: weekend letter, left edge in cm, 1 when it is the current weekend.
% The header band carries the weekend label on vellum; on the current plate the rest of the
% band is hatched.
\newcommand{\cxplate}[3]{%
  \begin{scope}[shift={(#2,\cxtop)}]
    \ifnum#3=1
      \fill[hatch] (0,0) rectangle (\cxW,\cxband);
      \def\cxlab{brand}\def\cxout{brand}\def\cxoutw{1.1pt}
    \else
      \def\cxlab{muted}\def\cxout{ruledark}\def\cxoutw{0.65pt}
    \fi
    \node[anchor=west, inner xsep=0.28cm, inner ysep=0, minimum height=\cxband cm, fill=paper,
      text=muted, font=\small\itshape] (lab#1) at (0,{\cxband/2}) {Weekend \csname wk#1Num\endcsname};
    \ifnum#3=1
      \draw[brand, line width=0.5pt] (lab#1.north east) -- (lab#1.south east);
    \fi
    \draw[\cxout, line width=\cxoutw, line join=miter] (0,0) rectangle (\cxW,\cxH);
    \draw[\cxout, line width=0.45pt] (0,\cxband) -- (\cxW,\cxband);
    \ifnum#3=1
      \node[anchor=south west, inner sep=0, fill=accent2, text=paper, font=\small\itshape,
        minimum height=0.46cm, inner xsep=2.6mm] at (0,0) {You are here};
    \fi
    \node[anchor=base west, inner sep=0, text=\cxlab, font=\fontsize{14}{16}\selectfont\itshape]
      at (0.28,{\cxband+0.54}) {\cxMuse{#1}};
    \node[anchor=north west, inner sep=0, text width=3.6cm, align=flush left,
      font=\fontsize{13}{14.4}\selectfont\bfseries] at (0.28,{\cxband+0.82}) {\csname wk#1Title\endcsname};
    \node[anchor=base west, inner sep=0, font=\itshape] at (0.28,{\cxband+2.19})
      {\csname wk#1Dates\endcsname};
    \draw[ruling, line width=0.4pt] (0.28,{\cxband+2.4}) -- (1.1,{\cxband+2.4});
    \node[anchor=north west, inner sep=0, text width=3.6cm, align=flush left,
      font=\fontsize{9}{10.4}\selectfont] at (0.28,{\cxband+2.52})
      {{\color{muted}\csname wk#1Topics\endcsname}\par\vskip 1.2mm
       {\itshape\color{muted}Project:} \csname wk#1Project\endcsname};
  \end{scope}}

\newcommand{\currentweekend}{a}   % the weekend that carries the hatched band and the flag
\newcommand{\threeweekends}{%
  \centering
  \def\cxcura{0}\def\cxcurb{0}\def\cxcurc{0}%
  \expandafter\def\csname cxcur\currentweekend\endcsname{1}%
  \pgfmathsetmacro{\cxxb}{\cxW+\cxG}
  \pgfmathsetmacro{\cxxc}{2*(\cxW+\cxG)}
  \pgfmathsetmacro{\cxxg}{2*\cxW+1.5*\cxG}
  \begin{tikzpicture}[x=1cm, y=-1cm]
    \path[use as bounding box] (0,0) rectangle (13,{\cxtop+\cxH});
    \cxplate{a}{0}{\cxcura}
    \cxplate{b}{\cxxb}{\cxcurb}
    \cxplate{c}{\cxxc}{\cxcurc}
    \ifcxpower
      \node[anchor=base, inner sep=0, text=muted, font=\itshape] (pow)
        at (\cxxg,{\cxtop-0.2}) {The Power};
      \draw[brand, line width=0.8pt] ([shift={(-0.8mm,0.19em)}]pow.base west) --
        ([shift={(0.8mm,0.19em)}]pow.base east);
      \draw[muted, line width=0.5pt, dash pattern=on 2.2pt off 1.8pt]
        (\cxxg,{\cxtop-0.1}) -- (\cxxg,{\cxtop+\cxH-0.1});
    \fi
  \end{tikzpicture}}
        % the facts of the course and the three weekends frame
%   \renewcommand{\deckslug}{AI in Industry, weekend 1}
% and is compiled with lualatex FROM the deck's own directory, three times: the leaf, the dots
% and the balloons are placed with remember picture and settle on the third pass. Fonts are
% EB Garamond, Garamond Math and Libertinus Mono (all in TeX Live) and Kefa III for Ethiopic
% (a macOS system font; set \geezfontpath before this file to use another folder).
%
% What a deck sets:
%   \title \subtitle \author \institute \date       the title page
%   \renewcommand{\deckslug}{...}                   the italic slug in the footer
%   \def\titlegeez{...}                             optional: a Ge'ez line above the cover title
%   \plateImage \plateCredit \quoteText \quoteCitation   the cold open, shown by \coldopenframe
% Frames 'title' and 'cold open' are [plain,noframenumbering]; a \section makes its own divider.
%
% The kit a deck may use:
%   macros   \hl{key words}  \geez{...}  \geezNumeral{n}  \cxballoon  \treelink  \weblink  \qrlink
%            \framedimgh{height}{file}  \coldopenframe  \begin{codeplate}...\end{codeplate}
%            \begin{cxlist}[\small]...\end{cxlist}  a denser list for reference slides
%   styles   chip, chip strong, chip dim, edge strong, edge dim, elab, elab dim, axis, vec a,
%            vec b, vlab a, vlab b, anglearc, anglelab, gloss, hatch, hatch fine
%   colours  paper ink muted brand accent2 good bad ruling ruledark nib
%
% House rules that bind every deck (the lecture-slides skill has the long form):
%   light and projection first; no red against green (good is blue, bad is vermillion); no ALL
%   CAPS; no dash or double hyphen in prose; nothing under a diagram; a section divider carries a
%   numeral, "Section N", the title and a band, nothing else; only the title page names the
%   lecturer; the only number in the chrome is the folio; no cell grids, column letters, tabs,
%   title blocks or arcade chrome. Never name a colour fg or bg: beamer rebinds both.

\usepackage{fontspec}
\usepackage{amsmath}
\usepackage{booktabs}
\usepackage{unicode-math}
\usefonttheme{professionalfonts}
\usefonttheme{serif}
\usepackage{tikz}
\usetikzlibrary{arrows.meta, positioning, calc, backgrounds, fit, overlay-beamer-styles,
                decorations.pathreplacing, patterns.meta}
\setbeamertemplate{navigation symbols}{}

% === type: EB Garamond (TeX Live), Garamond-Math for the maths, Kefa III for Ethiopic ===
\setmainfont{EBGaramond}[
  Extension      = .otf,
  UprightFont    = *-Medium,
  ItalicFont     = *-MediumItalic,
  BoldFont       = *-SemiBold,
  BoldItalicFont = *-SemiBoldItalic,
  Ligatures      = TeX,
  Numbers        = Lining,
  Scale          = 1.08]
\setmathfont{Garamond-Math.otf}[Scale=1.08]
% a monospaced face for code and addresses, sized to the lowercase of the main face
\setmonofont{LibertinusMono-Regular.otf}[Scale=MatchLowercase]

% Kefa III is a macOS system font. Its light and bold instances are variable font instances.
% Without the file, the deck still compiles: Ge'ez text then lacks glyphs and a warning says so.
\providecommand{\geezfontpath}{/System/Library/Fonts/Supplemental/}
\IfFileExists{\geezfontpath KefaIII.ttf}{%
  \newfontfamily\geezfont[Path=\geezfontpath]{KefaIII.ttf}
  % the light instance, for large numerals
  \newfontfamily\geezlightfont[Path=\geezfontpath, RawFeature={instance=Light}]{KefaIII.ttf}
  % the bold instance, for the numerals in balloons (the regular weight is too thin there)
  \newfontfamily\geezboldfont[Path=\geezfontpath, RawFeature={instance=Bold}]{KefaIII.ttf}
}{%
  \PackageWarningNoLine{ainit-preamble}{Kefa III was not found in \geezfontpath. Ge'ez text
    will lack glyphs. Set \string\geezfontpath\space before this file.}%
  \newcommand{\geezfont}{}\newcommand{\geezlightfont}{}\newcommand{\geezboldfont}{}}
\newcommand{\geez}[1]{{\geezfont #1}}
\newcommand{\geezlight}[1]{{\geezlightfont #1}}

% === colours. Never name one fg or bg (beamer rebinds both). ===
\definecolor{paper}{HTML}{F9F4E3}     % vellum, lightness above 93 percent
\definecolor{ink}{HTML}{241F1A}       % warm black, iron gall
\definecolor{muted}{HTML}{5E5448}     % secondary text, darker than black!55
\definecolor{brand}{HTML}{58409A}     % violet in the rubric role
\definecolor{accent2}{HTML}{8F5A06}   % ochre, the second accent
\definecolor{good}{HTML}{0072B2}      % Okabe-Ito blue: a win
\definecolor{bad}{HTML}{D55E00}       % Okabe-Ito vermillion: a loss. No red against green, ever.
\colorlet{ruling}{ink!40!paper}       % the ruled lines of the leaf
\colorlet{ruledark}{ink!62!paper}     % the heavier ruled lines
\colorlet{nib}{ink!88!paper}          % the outline ink of the figures

\setbeamercolor{normal text}{fg=ink, bg=paper}
\setbeamercolor{structure}{fg=brand}
\setbeamercolor{alerted text}{fg=brand}
\setbeamersize{text margin left=15mm, text margin right=15mm}

% === links: violet, and a raw address keeps the mono face, so a link is marked twice ===
% Beamer loads hyperref. The border is cleared; colorlinks stays off because it would also tint
% the section title on the divider.
\hypersetup{pdfborder={0 0 0}, breaklinks=true}
\urlstyle{tt}
\renewcommand{\UrlFont}{\ttfamily\color{brand}}
% \weblink{https://...}{the words the room reads}: an inline link whose visible text is prose.
\newcommand{\weblink}[2]{\href{#1}{\color{brand}#2}}

% \qrlink{2.6cm}{https://...}: a vector QR on a plain white plate in a hairline mount. The code is
% generated from the address at compile time, so it cannot disagree with the address printed
% beside it. Always print the address next to it.
% Two guards, both invisible in the source. nolink: qrcode makes the whole code a link by default,
% and a link is painted in the link colour. \color{black}: qrcode draws in the current colour, so
% on a full bleed frame where \color{white} is active the modules would vanish.
\usepackage{qrcode}
\newcommand{\qrlink}[2]{%
  \tikz{\node[inner sep=4pt, fill=white, draw=ruledark, line width=0.6pt, line join=miter]
    {\color{black}\qrcode[nolink, height=#1]{#2}};}}

% \framedimgh{2.5cm}{img/photo.jpg}: a photograph or figure of the given height in a hairline mount,
% standing 1.4 mm off the picture as the plate of the cold open does.
\newcommand{\framedimgh}[2]{%
  \tikz{\node[inner sep=1.4mm, draw=ruling, line width=0.4pt, line join=miter]
    {\includegraphics[height=#1]{#2}};}}

% === the overlay style used by every piece of page furniture ===
\tikzset{cx page/.style={remember picture, overlay, shift={(current page.north west)},
  x=1mm, y=-1mm}}

% === hatching: the one fill of the figures. Violet lines at 45 degrees, no solid colour under them ===
\tikzset{
  hatch/.style={pattern={Lines[angle=45, distance=0.95mm, line width=0.42pt]}, pattern color=brand},
  hatch fine/.style={pattern={Lines[angle=45, distance=0.62mm, line width=0.36pt]}, pattern color=brand},
}

% === Ge'ez numerals up to five, digits beyond ===
% The room learns ፩ to ፭ in the Ge'ez games (the course numbers the games with them).
\newcommand{\geezNumeral}[1]{\ifcase#1 0\or ፩\or ፪\or ፫\or ፬\or ፭\else\number#1\fi}
\newcommand{\cxnum}[1]{\ifnum#1<6\relax{\geezboldfont\geezNumeral{#1}}\else\number#1\fi}
\newcounter{glossn}

% a box that is scaled down to fit when its natural width exceeds #1, and left alone otherwise
\newsavebox{\cxfitbox}
\newcommand{\cxfit}[2]{%
  \sbox{\cxfitbox}{#2}%
  \ifdim\wd\cxfitbox>#1\relax\resizebox{#1}{!}{\usebox{\cxfitbox}}\else\usebox{\cxfitbox}\fi}

% === the footer slug: each deck renews it ===
\providecommand{\deckslug}{AI in Industry}

% === the leaf: ruled block with a fine double rule, pricking marks in the outer margin ===
\setbeamertemplate{background}{%
  \setcounter{glossn}{0}%
  \begin{tikzpicture}[cx page]
    \draw[ruling, line width=0.55pt] (8,5) rectangle (152,82);
    \draw[ruling, line width=0.25pt] (8.9,5.9) rectangle (151.1,81.1);
    \foreach \i in {0,...,12} {\fill[ruling] (3.2,{5+6.4167*\i}) circle (0.3);}
  \end{tikzpicture}}

% === the title sits above the first ruled line of the block; an over long title is scaled to fit ===
\setbeamertemplate{frametitle}{%
  \vbox to 17.5mm{\offinterlineskip
    \vskip 8.3mm
    \hbox to\textwidth{\vrule height 5.3mm depth 1.9mm width 0pt
      \cxfit{\textwidth}{\color{ink}\fontsize{19}{22}\selectfont\bfseries\insertframetitle}\hfil}%
    \vskip 0.8mm
    \hbox to\textwidth{\color{ruledark}\vrule height 0.7pt depth 0pt width\textwidth}%
    \vskip 0.6mm
    \hbox to\textwidth{\color{ruling}\vrule height 0.3pt depth 0pt width\textwidth}%
    \vss}}

% === the folio in the lower margin, the italic slug beside it ===
\setbeamertemplate{footline}{%
  \vbox to 8.2mm{\offinterlineskip\vss
    \hbox to\paperwidth{\hskip 15mm
      {\color{brand}\fontsize{12}{14}\selectfont\insertframenumber}\hskip 3.4mm
      {\color{muted}\footnotesize\itshape\deckslug}\hfil}%
    \vskip 2.2mm}}

% === the title page ===
% \titlegeez (a Ge'ez line over the title) and \geezYear (the year, in Ge'ez numerals, after the
% date) are optional. \geezYear comes from ainit-course.tex.
\setbeamertemplate{title page}{%
  \begin{tikzpicture}[cx page]
    \node[text=brand, font=\fontsize{17}{20}\selectfont, anchor=base] at (80,20.5)
      {\ifdefined\titlegeez\geez{\titlegeez}\fi};
    \node[text=ink, font=\fontsize{27}{30}\selectfont\bfseries, text width=104mm, align=center,
      anchor=north, inner sep=0] at (80,25) {\inserttitle};
    \draw[ruledark, line width=0.8pt] (62,50.5) -- (98,50.5);
    \draw[ruling, line width=0.3pt] (62,51.5) -- (98,51.5);
    \node[text=muted, font=\large\itshape, anchor=north, inner sep=0] at (80,55)
      {\insertsubtitle};
    \node[text=ink, font=\large, anchor=base, inner sep=0] at (80,67) {\insertauthor};
    \node[text=muted, font=\normalsize\itshape, anchor=base, inner sep=0] at (80,72.2)
      {\insertinstitute};
    \node[text=muted, font=\normalsize\itshape, anchor=base, inner sep=0] at (80,77.2)
      {\insertdate\ifdefined\geezYear\hspace{3.2mm}{\upshape\color{brand}\geez{\geezYear}}\fi};
  \end{tikzpicture}}

% === the section page: a numeral in a heavy square, Section N, the title, a hatched band ===
% The numeral is Ge'ez up to five and a digit beyond. A title that is wider than 92 mm at the
% display size wraps over several lines at a smaller size, and the band moves down under it.
\newif\ifcxwrap
\setbeamertemplate{section page}{%
  \sbox{\cxfitbox}{\fontsize{36}{40}\selectfont\bfseries\insertsectionhead}%
  \ifdim\wd\cxfitbox>92mm\relax\cxwraptrue\else\cxwrapfalse\fi
  \begin{tikzpicture}[cx page]
    \draw[brand, line width=1.8pt, line join=miter] (17,30.2) rectangle (44,57.2);
    \draw[brand, line width=0.4pt] (18.9,32.1) rectangle (42.1,55.3);
    \ifnum\insertsectionnumber<6\relax
      \node[text=brand, font=\fontsize{54}{54}\selectfont, anchor=center] at (30.56,43.7)
        {\geezlight{\geezNumeral{\insertsectionnumber}}};
    \else
      \node[text=brand, font=\fontsize{54}{54}\selectfont, anchor=center] at (30.56,43.7)
        {\number\insertsectionnumber};
    \fi
    \node[text=muted, font=\fontsize{14}{16}\selectfont\itshape, anchor=base west, inner sep=0]
      at (53,33.2) {Section \insertsectionnumber};
    \ifcxwrap
      \node[text=ink, font=\fontsize{26}{29}\selectfont\bfseries, anchor=north west, inner sep=0,
        text width=92mm, align=flush left] (cxst) at (53,36.5) {\insertsectionhead};
      \fill[hatch] ([shift={(0,3.4)}]cxst.south west) rectangle ([shift={(92,7)}]cxst.south west);
      \draw[brand, line width=0.5pt] ([shift={(0,3.4)}]cxst.south west) -- ++(92,0)
        ([shift={(0,7)}]cxst.south west) -- ++(92,0);
    \else
      \node[text=ink, font=\fontsize{36}{40}\selectfont\bfseries, anchor=base west, inner sep=0]
        at (53,47.5) {\insertsectionhead};
      \fill[hatch] (53,53.6) rectangle (145,57.2);
      \draw[brand, line width=0.5pt] (53,53.6) -- (145,53.6) (53,57.2) -- (145,57.2);
    \fi
  \end{tikzpicture}}
\AtBeginSection[]{\begin{frame}[plain,noframenumbering]\usebeamertemplate{section page}\end{frame}}

% === rubrication: key words in semibold violet ===
\newcommand{\hl}[1]{\textcolor{brand}{\bfseries #1}}

% === the cold open: the plate in a hairline mount, the verse, the citation ===
% The deck sets \plateImage (a file name), \plateCredit, \quoteText (lines separated by \\, key
% words in \hl{}) and \quoteCitation, and calls \coldopenframe inside a [plain] frame.
\newcommand{\coldopenframe}{%
  \begin{tikzpicture}[cx page]
    \node[anchor=north west, inner sep=0] (plate) at (16,15)
      {\includegraphics[height=52mm]{\plateImage}};
    \draw[ruling, line width=0.4pt] ([shift={(-1.6,-1.6)}]plate.north west) rectangle
      ([shift={(1.6,1.6)}]plate.south east);
    \node[anchor=north, inner sep=0, text=muted, font=\scriptsize\itshape]
      at ($(plate.south)+(0,3.2)$) {\plateCredit};
    \node[anchor=west, inner sep=0, text width=86mm, align=flush left,
      font=\fontsize{11.5}{17}\selectfont] at (62,43.5) {%
      \quoteText\par\vskip 4mm
      {\color{ruling}\vrule height 0.4pt depth 0pt width 22mm}\par\vskip 2mm
      {\small\itshape\color{muted}\quoteCitation}};
  \end{tikzpicture}}

% === a numbered balloon: a hairline circle with a numeral, centre (#1,#2) in the picture's units ===
\newcommand{\cxballoon}[4]{% x, y, radius in mm, numeral
  \draw[brand, line width=0.7pt, fill=paper] (#1,#2) circle [radius=#3mm];
  \node[text=brand, inner sep=0, font=\fontsize{14}{14}\selectfont] at (#1,#2) {\cxnum{#4}};}

% \treelink{edge style}{label style}{from}{to}{label}: one right angled branch of a tree.
% It leaves the east anchor of the parent, runs level to a junction at two fifths of the gap,
% turns once to the height of the child and arrives level at its west anchor. The first branch
% drawn from a parent carries the shared stretch, so declare the chosen branch first: later
% branches from the same parent begin at the junction, and a dashed branch never overprints the
% heavier path.
\newcommand{\treelink}[5]{%
  \ifcsname cx@trunk@#3\endcsname
    \draw[#1] let \p1=(#3.east), \p2=(#4.west), \p3=($(\p1)!0.4!(\p2)$),
                  \p4=($(\x3,\y2)!0.5!(\p2)$) in
      (\p4) node[#2, above=0.4mm] {#5}
      (\x3,\y1) -- (\x3,\y2) -- (\p2);
  \else
    \expandafter\def\csname cx@trunk@#3\endcsname{}%
    \draw[#1] let \p1=(#3.east), \p2=(#4.west), \p3=($(\p1)!0.4!(\p2)$),
                  \p4=($(\x3,\y2)!0.5!(\p2)$) in
      (\p4) node[#2, above=0.4mm] {#5}
      (\p1) -- (\x3,\y1) -- (\x3,\y2) -- (\p2);
  \fi}

% === list markers: a small square hatched in violet ===
\newcommand{\cxhatchsq}[1]{%
  \raisebox{0.12ex}{\tikz[baseline=0pt]{%
    \fill[hatch fine] (0,0) rectangle (#1,#1);
    \draw[brand, line width=0.55pt, line join=miter] (0,0) rectangle (#1,#1);}}}
\setbeamertemplate{itemize item}{\cxhatchsq{3mm}}
\setbeamertemplate{itemize subitem}{\cxhatchsq{2.1mm}}
\setbeamertemplate{enumerate item}{%
  \tikz[baseline=(enumn.base)]\node[draw=brand, line width=0.7pt, circle, inner sep=0.5pt,
    minimum size=5.4mm, font=\small, text=brand, fill=paper] (enumn) {\cxnum{\insertenumlabel}};}
\setbeamertemplate{enumerate subitem}{\insertenumlabel.}
\setbeamerfont{itemize/enumerate body}{size=\Large}
\setbeamerfont{itemize/enumerate subbody}{size=\large}
\setbeamertemplate{itemize/enumerate body begin}{\vspace*{0mm}}
% beamer resets the list parameters inside \list, so the spacing is set in the list defaults
\makeatletter
\def\@listi{\leftmargin\leftmargini
            \topsep 0\p@ \parsep 0\p@
            \itemsep 6mm \@plus1mm \@minus1mm}
\let\@listI\@listi
\def\@listii{\leftmargin\leftmarginii
             \topsep 0\p@ \parsep 0\p@
             \itemsep 2mm \@plus0.5mm \@minus0.5mm}
\makeatother
% a denser list for reference slides (administration, rosters, small print): a smaller marker and
% the text size #1, \small by default. Set \itemsep right after \begin{cxlist} as with itemize.
\newenvironment{cxlist}[1][\small]{%
  \begingroup
  \setbeamertemplate{itemize item}{\cxhatchsq{2.2mm}}%
  \setbeamerfont{itemize/enumerate body}{size=#1}%
  \begin{itemize}}{\end{itemize}\endgroup}

% === the code plate: mono text in a hairline box under a hatched band ===
% Use it in a [fragile] frame. It is built for lines up to about 60 characters: a longer line
% runs past the box, so shorten the code or the prompt instead of shrinking the type.
\usepackage{listings}
\usepackage[skins,listings]{tcolorbox}
\lstdefinestyle{ainit}{
  basicstyle=\ttfamily\normalsize, keywordstyle=\color{brand}\bfseries,
  commentstyle=\color{muted}\itshape, stringstyle=\color{accent2},
  showstringspaces=false, columns=fullflexible, keepspaces=true, breaklines=false,
  aboveskip=0pt, belowskip=0pt, language=Python}
\newtcblisting{codeplate}[1][]{%
  enhanced, listing only, listing options={style=ainit},
  colback=paper, colframe=ruledark, boxrule=0.6pt, sharp corners,
  left=4mm, right=4mm, top=5.6mm, bottom=2.2mm, before skip=1mm, after skip=1mm,
  overlay unbroken={%
    \fill[hatch] ([yshift=-2.9mm]frame.north west) rectangle (frame.north east);
    \draw[brand, line width=0.5pt]
      ([yshift=-2.9mm]frame.north west) -- ([yshift=-2.9mm]frame.north east);},
  #1}

% === the TikZ styles of the figures ===
\tikzset{
  >={Triangle[length=2.5mm, width=1.9mm]},
  chip/.style={draw=nib, line width=0.7pt, line join=miter, minimum width=1.85cm,
    minimum height=0.82cm, inner sep=3pt, text=ink, fill=paper, font=\large},
  chip strong/.style={chip, draw=brand, line width=1.0pt,
    path picture={%
      \fill[hatch, even odd rule]
        (path picture bounding box.south west) rectangle (path picture bounding box.north east)
        ([shift={(1.25mm,1.25mm)}]path picture bounding box.south west) rectangle
        ([shift={(-1.25mm,-1.25mm)}]path picture bounding box.north east);
      \draw[brand, line width=0.35pt]
        ([shift={(1.25mm,1.25mm)}]path picture bounding box.south west) rectangle
        ([shift={(-1.25mm,-1.25mm)}]path picture bounding box.north east);}},
  chip dim/.style={chip, draw=muted, line width=0.6pt, dash pattern=on 2.4pt off 1.8pt,
    text=muted},
  edge strong/.style={draw=brand, line width=1.3pt, line join=miter,
    -{Triangle[length=2.9mm, width=2.2mm]}},
  edge dim/.style={draw=muted, line width=0.6pt, line join=miter, dash pattern=on 3pt off 2.2pt,
    -{Triangle[length=2.5mm, width=1.8mm]}},
  elab/.style={font=\small\itshape, text=brand, fill=paper, inner sep=1.5pt},
  elab dim/.style={font=\small\itshape, text=muted, fill=paper, inner sep=1.5pt},
  axis/.style={draw=muted, line width=0.6pt, -{Triangle[length=2.3mm, width=1.6mm]}},
  vec a/.style={draw=brand, line width=1.4pt, -{Triangle[length=3.1mm, width=2.3mm]}},
  vec b/.style={draw=accent2, line width=1.4pt, -{Triangle[length=3.1mm, width=2.3mm]}},
  vlab a/.style={text=brand, font=\normalsize},
  vlab b/.style={text=accent2, font=\normalsize},
  anglearc/.style={draw=ink, line width=0.6pt,
    {Triangle[length=2mm, width=1.4mm]}-{Triangle[length=2mm, width=1.4mm]}},
  anglelab/.style={text=ink, font=\normalsize},
  % a note: a numbered balloon on a hairline leader, with its text. The balloons number
  % themselves 1, 2, 3 on each page.
  gloss/.style={font=\small\itshape, text=muted, text width=3.9cm, align=flush left,
    inner xsep=6.9mm, inner ysep=0.6mm,
    execute at begin node={\stepcounter{glossn}},
    path picture={%
      \draw[brand, line width=0.5pt]
        ($(path picture bounding box.north west)+(2.9mm,-5.4mm)$) --
        ($(path picture bounding box.south west)+(2.9mm,0.3mm)$);
      \draw[brand, line width=0.7pt, fill=paper]
        ($(path picture bounding box.north west)+(2.9mm,-2.9mm)$) circle (2.5mm);
      \pgftext[at={\pgfpointadd{\pgfpointanchor{path picture bounding box}{north west}}%
        {\pgfpoint{2.9mm}{-2.9mm}}}]
        {\color{brand}\fontsize{10.5}{10.5}\selectfont\upshape\cxnum{\value{glossn}}}}},
}
      % this file
%   % Copied on 2026-11-03 from eth-ainit-hs26/w1-lecture-material, commit 9792819, file style/ainit-course.tex. To pull a fix, diff against the source and recopy.
% ===== ainit-course.tex: the facts of AINIT HS26 and the three weekends frame =====
% Load this after ainit-preamble.tex in any deck that shows the course roadmap or the title page's
% Ge'ez year. The facts live here once, so a change of date or title is one edit.
%
% Where the facts come from. Dates, themes and titles are the master spreadsheet's Syllabus tab and
% agree with the course website. The three project names differ between the website, the sheet and
% the introduction deck. This file follows the introduction deck (admin-deck/), because that is the
% deck the participants see first. The topics line of each weekend is a short reading of the
% weekend's slots. Re-read the sheet before the semester starts and after any change of schedule.
%
% For the next academic year change \geezYear, the dates, and whatever else moved.

% 2026 in Ge'ez numerals: 20 x 100 + 20 + 6. The title page prints it after the date.
\newcommand{\geezYear}{፳፻፳፮}

% ---- the three weekends (a, b, c): number, title, dates, topics, project, Faust epithet ----
% The epithets come from the course's cold open: Goethe's Faust renders the Greek word logos as
% the Word, the Thought and the Act, and refuses the Power. They map to the three weekends.
\newcommand{\wkaNum}{1}\newcommand{\wkaTitle}{Large language models}
\newcommand{\wkaDates}{16 and 17 October}
\newcommand{\wkaTopics}{Ge'ez games, prompting, OPRO, MLOps}
\newcommand{\wkaProject}{LLM-based marketing agent}\newcommand{\wkaMuse}{the Word}
\newcommand{\wkbNum}{2}\newcommand{\wkbTitle}{Representation learning}
\newcommand{\wkbDates}{13 and 14 November}
\newcommand{\wkbTopics}{Recommendation systems, matrix factorization, bandits}
\newcommand{\wkbProject}{Recommendation systems}\newcommand{\wkbMuse}{the Thought}
\newcommand{\wkcNum}{3}\newcommand{\wkcTitle}{Reinforcement learning}
\newcommand{\wkcDates}{11 and 12 December}
\newcommand{\wkcTopics}{MDPs, Bellman equations, SARSA, REINFORCE}
\newcommand{\wkcProject}{Reinforcement learning for delivery optimization}
\newcommand{\wkcMuse}{the Act}

% ===== frame: the three weekends, three plates side by side =====
% The current weekend has a hatched header band and an ochre flag. Set it with
% \renewcommand{\currentweekend}{b} in a weekend 2 deck. The refused Power hangs as a struck label
% above the gap between the second and the third plate, with a dashed phantom line running down that
% gap. Switch it off with \cxpowerfalse after loading this file.
\newif\ifcxpower
\cxpowertrue
\newcommand{\cxW}{4.08}   % plate width in cm
\newcommand{\cxG}{0.38}   % gap between plates in cm
\newcommand{\cxH}{5.0}    % plate height in cm
\newcommand{\cxtop}{0.56} % height of the row of flags above the plates, in cm
\newcommand{\cxband}{0.44}% height of the header band of a plate, in cm

% the capital form of a Muse: "the Word" becomes "The Word"
\def\cxcaptheA the #1\relax{The #1}
\newcommand{\cxMuse}[1]{{\edef\cxtmp{\csname wk#1Muse\endcsname}\expandafter\cxcaptheA\cxtmp\relax}}

% one plate: weekend letter, left edge in cm, 1 when it is the current weekend.
% The header band carries the weekend label on vellum; on the current plate the rest of the
% band is hatched.
\newcommand{\cxplate}[3]{%
  \begin{scope}[shift={(#2,\cxtop)}]
    \ifnum#3=1
      \fill[hatch] (0,0) rectangle (\cxW,\cxband);
      \def\cxlab{brand}\def\cxout{brand}\def\cxoutw{1.1pt}
    \else
      \def\cxlab{muted}\def\cxout{ruledark}\def\cxoutw{0.65pt}
    \fi
    \node[anchor=west, inner xsep=0.28cm, inner ysep=0, minimum height=\cxband cm, fill=paper,
      text=muted, font=\small\itshape] (lab#1) at (0,{\cxband/2}) {Weekend \csname wk#1Num\endcsname};
    \ifnum#3=1
      \draw[brand, line width=0.5pt] (lab#1.north east) -- (lab#1.south east);
    \fi
    \draw[\cxout, line width=\cxoutw, line join=miter] (0,0) rectangle (\cxW,\cxH);
    \draw[\cxout, line width=0.45pt] (0,\cxband) -- (\cxW,\cxband);
    \ifnum#3=1
      \node[anchor=south west, inner sep=0, fill=accent2, text=paper, font=\small\itshape,
        minimum height=0.46cm, inner xsep=2.6mm] at (0,0) {You are here};
    \fi
    \node[anchor=base west, inner sep=0, text=\cxlab, font=\fontsize{14}{16}\selectfont\itshape]
      at (0.28,{\cxband+0.54}) {\cxMuse{#1}};
    \node[anchor=north west, inner sep=0, text width=3.6cm, align=flush left,
      font=\fontsize{13}{14.4}\selectfont\bfseries] at (0.28,{\cxband+0.82}) {\csname wk#1Title\endcsname};
    \node[anchor=base west, inner sep=0, font=\itshape] at (0.28,{\cxband+2.19})
      {\csname wk#1Dates\endcsname};
    \draw[ruling, line width=0.4pt] (0.28,{\cxband+2.4}) -- (1.1,{\cxband+2.4});
    \node[anchor=north west, inner sep=0, text width=3.6cm, align=flush left,
      font=\fontsize{9}{10.4}\selectfont] at (0.28,{\cxband+2.52})
      {{\color{muted}\csname wk#1Topics\endcsname}\par\vskip 1.2mm
       {\itshape\color{muted}Project:} \csname wk#1Project\endcsname};
  \end{scope}}

\newcommand{\currentweekend}{a}   % the weekend that carries the hatched band and the flag
\newcommand{\threeweekends}{%
  \centering
  \def\cxcura{0}\def\cxcurb{0}\def\cxcurc{0}%
  \expandafter\def\csname cxcur\currentweekend\endcsname{1}%
  \pgfmathsetmacro{\cxxb}{\cxW+\cxG}
  \pgfmathsetmacro{\cxxc}{2*(\cxW+\cxG)}
  \pgfmathsetmacro{\cxxg}{2*\cxW+1.5*\cxG}
  \begin{tikzpicture}[x=1cm, y=-1cm]
    \path[use as bounding box] (0,0) rectangle (13,{\cxtop+\cxH});
    \cxplate{a}{0}{\cxcura}
    \cxplate{b}{\cxxb}{\cxcurb}
    \cxplate{c}{\cxxc}{\cxcurc}
    \ifcxpower
      \node[anchor=base, inner sep=0, text=muted, font=\itshape] (pow)
        at (\cxxg,{\cxtop-0.2}) {The Power};
      \draw[brand, line width=0.8pt] ([shift={(-0.8mm,0.19em)}]pow.base west) --
        ([shift={(0.8mm,0.19em)}]pow.base east);
      \draw[muted, line width=0.5pt, dash pattern=on 2.2pt off 1.8pt]
        (\cxxg,{\cxtop-0.1}) -- (\cxxg,{\cxtop+\cxH-0.1});
    \fi
  \end{tikzpicture}}
        % the facts of the course and the three weekends frame
%   \renewcommand{\deckslug}{AI in Industry, weekend 1}
% and is compiled with lualatex FROM the deck's own directory, three times: the leaf, the dots
% and the balloons are placed with remember picture and settle on the third pass. Fonts are
% EB Garamond, Garamond Math and Libertinus Mono (all in TeX Live) and Kefa III for Ethiopic
% (a macOS system font; set \geezfontpath before this file to use another folder).
%
% What a deck sets:
%   \title \subtitle \author \institute \date       the title page
%   \renewcommand{\deckslug}{...}                   the italic slug in the footer
%   \def\titlegeez{...}                             optional: a Ge'ez line above the cover title
%   \plateImage \plateCredit \quoteText \quoteCitation   the cold open, shown by \coldopenframe
% Frames 'title' and 'cold open' are [plain,noframenumbering]; a \section makes its own divider.
%
% The kit a deck may use:
%   macros   \hl{key words}  \geez{...}  \geezNumeral{n}  \cxballoon  \treelink  \weblink  \qrlink
%            \framedimgh{height}{file}  \coldopenframe  \begin{codeplate}...\end{codeplate}
%            \begin{cxlist}[\small]...\end{cxlist}  a denser list for reference slides
%   styles   chip, chip strong, chip dim, edge strong, edge dim, elab, elab dim, axis, vec a,
%            vec b, vlab a, vlab b, anglearc, anglelab, gloss, hatch, hatch fine
%   colours  paper ink muted brand accent2 good bad ruling ruledark nib
%
% House rules that bind every deck (the lecture-slides skill has the long form):
%   light and projection first; no red against green (good is blue, bad is vermillion); no ALL
%   CAPS; no dash or double hyphen in prose; nothing under a diagram; a section divider carries a
%   numeral, "Section N", the title and a band, nothing else; only the title page names the
%   lecturer; the only number in the chrome is the folio; no cell grids, column letters, tabs,
%   title blocks or arcade chrome. Never name a colour fg or bg: beamer rebinds both.

\usepackage{fontspec}
\usepackage{amsmath}
\usepackage{booktabs}
\usepackage{unicode-math}
\usefonttheme{professionalfonts}
\usefonttheme{serif}
\usepackage{tikz}
\usetikzlibrary{arrows.meta, positioning, calc, backgrounds, fit, overlay-beamer-styles,
                decorations.pathreplacing, patterns.meta}
\setbeamertemplate{navigation symbols}{}

% === type: EB Garamond (TeX Live), Garamond-Math for the maths, Kefa III for Ethiopic ===
\setmainfont{EBGaramond}[
  Extension      = .otf,
  UprightFont    = *-Medium,
  ItalicFont     = *-MediumItalic,
  BoldFont       = *-SemiBold,
  BoldItalicFont = *-SemiBoldItalic,
  Ligatures      = TeX,
  Numbers        = Lining,
  Scale          = 1.08]
\setmathfont{Garamond-Math.otf}[Scale=1.08]
% a monospaced face for code and addresses, sized to the lowercase of the main face
\setmonofont{LibertinusMono-Regular.otf}[Scale=MatchLowercase]

% Kefa III is a macOS system font. Its light and bold instances are variable font instances.
% Without the file, the deck still compiles: Ge'ez text then lacks glyphs and a warning says so.
\providecommand{\geezfontpath}{/System/Library/Fonts/Supplemental/}
\IfFileExists{\geezfontpath KefaIII.ttf}{%
  \newfontfamily\geezfont[Path=\geezfontpath]{KefaIII.ttf}
  % the light instance, for large numerals
  \newfontfamily\geezlightfont[Path=\geezfontpath, RawFeature={instance=Light}]{KefaIII.ttf}
  % the bold instance, for the numerals in balloons (the regular weight is too thin there)
  \newfontfamily\geezboldfont[Path=\geezfontpath, RawFeature={instance=Bold}]{KefaIII.ttf}
}{%
  \PackageWarningNoLine{ainit-preamble}{Kefa III was not found in \geezfontpath. Ge'ez text
    will lack glyphs. Set \string\geezfontpath\space before this file.}%
  \newcommand{\geezfont}{}\newcommand{\geezlightfont}{}\newcommand{\geezboldfont}{}}
\newcommand{\geez}[1]{{\geezfont #1}}
\newcommand{\geezlight}[1]{{\geezlightfont #1}}

% === colours. Never name one fg or bg (beamer rebinds both). ===
\definecolor{paper}{HTML}{F9F4E3}     % vellum, lightness above 93 percent
\definecolor{ink}{HTML}{241F1A}       % warm black, iron gall
\definecolor{muted}{HTML}{5E5448}     % secondary text, darker than black!55
\definecolor{brand}{HTML}{58409A}     % violet in the rubric role
\definecolor{accent2}{HTML}{8F5A06}   % ochre, the second accent
\definecolor{good}{HTML}{0072B2}      % Okabe-Ito blue: a win
\definecolor{bad}{HTML}{D55E00}       % Okabe-Ito vermillion: a loss. No red against green, ever.
\colorlet{ruling}{ink!40!paper}       % the ruled lines of the leaf
\colorlet{ruledark}{ink!62!paper}     % the heavier ruled lines
\colorlet{nib}{ink!88!paper}          % the outline ink of the figures

\setbeamercolor{normal text}{fg=ink, bg=paper}
\setbeamercolor{structure}{fg=brand}
\setbeamercolor{alerted text}{fg=brand}
\setbeamersize{text margin left=15mm, text margin right=15mm}

% === links: violet, and a raw address keeps the mono face, so a link is marked twice ===
% Beamer loads hyperref. The border is cleared; colorlinks stays off because it would also tint
% the section title on the divider.
\hypersetup{pdfborder={0 0 0}, breaklinks=true}
\urlstyle{tt}
\renewcommand{\UrlFont}{\ttfamily\color{brand}}
% \weblink{https://...}{the words the room reads}: an inline link whose visible text is prose.
\newcommand{\weblink}[2]{\href{#1}{\color{brand}#2}}

% \qrlink{2.6cm}{https://...}: a vector QR on a plain white plate in a hairline mount. The code is
% generated from the address at compile time, so it cannot disagree with the address printed
% beside it. Always print the address next to it.
% Two guards, both invisible in the source. nolink: qrcode makes the whole code a link by default,
% and a link is painted in the link colour. \color{black}: qrcode draws in the current colour, so
% on a full bleed frame where \color{white} is active the modules would vanish.
\usepackage{qrcode}
\newcommand{\qrlink}[2]{%
  \tikz{\node[inner sep=4pt, fill=white, draw=ruledark, line width=0.6pt, line join=miter]
    {\color{black}\qrcode[nolink, height=#1]{#2}};}}

% \framedimgh{2.5cm}{img/photo.jpg}: a photograph or figure of the given height in a hairline mount,
% standing 1.4 mm off the picture as the plate of the cold open does.
\newcommand{\framedimgh}[2]{%
  \tikz{\node[inner sep=1.4mm, draw=ruling, line width=0.4pt, line join=miter]
    {\includegraphics[height=#1]{#2}};}}

% === the overlay style used by every piece of page furniture ===
\tikzset{cx page/.style={remember picture, overlay, shift={(current page.north west)},
  x=1mm, y=-1mm}}

% === hatching: the one fill of the figures. Violet lines at 45 degrees, no solid colour under them ===
\tikzset{
  hatch/.style={pattern={Lines[angle=45, distance=0.95mm, line width=0.42pt]}, pattern color=brand},
  hatch fine/.style={pattern={Lines[angle=45, distance=0.62mm, line width=0.36pt]}, pattern color=brand},
}

% === Ge'ez numerals up to five, digits beyond ===
% The room learns ፩ to ፭ in the Ge'ez games (the course numbers the games with them).
\newcommand{\geezNumeral}[1]{\ifcase#1 0\or ፩\or ፪\or ፫\or ፬\or ፭\else\number#1\fi}
\newcommand{\cxnum}[1]{\ifnum#1<6\relax{\geezboldfont\geezNumeral{#1}}\else\number#1\fi}
\newcounter{glossn}

% a box that is scaled down to fit when its natural width exceeds #1, and left alone otherwise
\newsavebox{\cxfitbox}
\newcommand{\cxfit}[2]{%
  \sbox{\cxfitbox}{#2}%
  \ifdim\wd\cxfitbox>#1\relax\resizebox{#1}{!}{\usebox{\cxfitbox}}\else\usebox{\cxfitbox}\fi}

% === the footer slug: each deck renews it ===
\providecommand{\deckslug}{AI in Industry}

% === the leaf: ruled block with a fine double rule, pricking marks in the outer margin ===
\setbeamertemplate{background}{%
  \setcounter{glossn}{0}%
  \begin{tikzpicture}[cx page]
    \draw[ruling, line width=0.55pt] (8,5) rectangle (152,82);
    \draw[ruling, line width=0.25pt] (8.9,5.9) rectangle (151.1,81.1);
    \foreach \i in {0,...,12} {\fill[ruling] (3.2,{5+6.4167*\i}) circle (0.3);}
  \end{tikzpicture}}

% === the title sits above the first ruled line of the block; an over long title is scaled to fit ===
\setbeamertemplate{frametitle}{%
  \vbox to 17.5mm{\offinterlineskip
    \vskip 8.3mm
    \hbox to\textwidth{\vrule height 5.3mm depth 1.9mm width 0pt
      \cxfit{\textwidth}{\color{ink}\fontsize{19}{22}\selectfont\bfseries\insertframetitle}\hfil}%
    \vskip 0.8mm
    \hbox to\textwidth{\color{ruledark}\vrule height 0.7pt depth 0pt width\textwidth}%
    \vskip 0.6mm
    \hbox to\textwidth{\color{ruling}\vrule height 0.3pt depth 0pt width\textwidth}%
    \vss}}

% === the folio in the lower margin, the italic slug beside it ===
\setbeamertemplate{footline}{%
  \vbox to 8.2mm{\offinterlineskip\vss
    \hbox to\paperwidth{\hskip 15mm
      {\color{brand}\fontsize{12}{14}\selectfont\insertframenumber}\hskip 3.4mm
      {\color{muted}\footnotesize\itshape\deckslug}\hfil}%
    \vskip 2.2mm}}

% === the title page ===
% \titlegeez (a Ge'ez line over the title) and \geezYear (the year, in Ge'ez numerals, after the
% date) are optional. \geezYear comes from ainit-course.tex.
\setbeamertemplate{title page}{%
  \begin{tikzpicture}[cx page]
    \node[text=brand, font=\fontsize{17}{20}\selectfont, anchor=base] at (80,20.5)
      {\ifdefined\titlegeez\geez{\titlegeez}\fi};
    \node[text=ink, font=\fontsize{27}{30}\selectfont\bfseries, text width=104mm, align=center,
      anchor=north, inner sep=0] at (80,25) {\inserttitle};
    \draw[ruledark, line width=0.8pt] (62,50.5) -- (98,50.5);
    \draw[ruling, line width=0.3pt] (62,51.5) -- (98,51.5);
    \node[text=muted, font=\large\itshape, anchor=north, inner sep=0] at (80,55)
      {\insertsubtitle};
    \node[text=ink, font=\large, anchor=base, inner sep=0] at (80,67) {\insertauthor};
    \node[text=muted, font=\normalsize\itshape, anchor=base, inner sep=0] at (80,72.2)
      {\insertinstitute};
    \node[text=muted, font=\normalsize\itshape, anchor=base, inner sep=0] at (80,77.2)
      {\insertdate\ifdefined\geezYear\hspace{3.2mm}{\upshape\color{brand}\geez{\geezYear}}\fi};
  \end{tikzpicture}}

% === the section page: a numeral in a heavy square, Section N, the title, a hatched band ===
% The numeral is Ge'ez up to five and a digit beyond. A title that is wider than 92 mm at the
% display size wraps over several lines at a smaller size, and the band moves down under it.
\newif\ifcxwrap
\setbeamertemplate{section page}{%
  \sbox{\cxfitbox}{\fontsize{36}{40}\selectfont\bfseries\insertsectionhead}%
  \ifdim\wd\cxfitbox>92mm\relax\cxwraptrue\else\cxwrapfalse\fi
  \begin{tikzpicture}[cx page]
    \draw[brand, line width=1.8pt, line join=miter] (17,30.2) rectangle (44,57.2);
    \draw[brand, line width=0.4pt] (18.9,32.1) rectangle (42.1,55.3);
    \ifnum\insertsectionnumber<6\relax
      \node[text=brand, font=\fontsize{54}{54}\selectfont, anchor=center] at (30.56,43.7)
        {\geezlight{\geezNumeral{\insertsectionnumber}}};
    \else
      \node[text=brand, font=\fontsize{54}{54}\selectfont, anchor=center] at (30.56,43.7)
        {\number\insertsectionnumber};
    \fi
    \node[text=muted, font=\fontsize{14}{16}\selectfont\itshape, anchor=base west, inner sep=0]
      at (53,33.2) {Section \insertsectionnumber};
    \ifcxwrap
      \node[text=ink, font=\fontsize{26}{29}\selectfont\bfseries, anchor=north west, inner sep=0,
        text width=92mm, align=flush left] (cxst) at (53,36.5) {\insertsectionhead};
      \fill[hatch] ([shift={(0,3.4)}]cxst.south west) rectangle ([shift={(92,7)}]cxst.south west);
      \draw[brand, line width=0.5pt] ([shift={(0,3.4)}]cxst.south west) -- ++(92,0)
        ([shift={(0,7)}]cxst.south west) -- ++(92,0);
    \else
      \node[text=ink, font=\fontsize{36}{40}\selectfont\bfseries, anchor=base west, inner sep=0]
        at (53,47.5) {\insertsectionhead};
      \fill[hatch] (53,53.6) rectangle (145,57.2);
      \draw[brand, line width=0.5pt] (53,53.6) -- (145,53.6) (53,57.2) -- (145,57.2);
    \fi
  \end{tikzpicture}}
\AtBeginSection[]{\begin{frame}[plain,noframenumbering]\usebeamertemplate{section page}\end{frame}}

% === rubrication: key words in semibold violet ===
\newcommand{\hl}[1]{\textcolor{brand}{\bfseries #1}}

% === the cold open: the plate in a hairline mount, the verse, the citation ===
% The deck sets \plateImage (a file name), \plateCredit, \quoteText (lines separated by \\, key
% words in \hl{}) and \quoteCitation, and calls \coldopenframe inside a [plain] frame.
\newcommand{\coldopenframe}{%
  \begin{tikzpicture}[cx page]
    \node[anchor=north west, inner sep=0] (plate) at (16,15)
      {\includegraphics[height=52mm]{\plateImage}};
    \draw[ruling, line width=0.4pt] ([shift={(-1.6,-1.6)}]plate.north west) rectangle
      ([shift={(1.6,1.6)}]plate.south east);
    \node[anchor=north, inner sep=0, text=muted, font=\scriptsize\itshape]
      at ($(plate.south)+(0,3.2)$) {\plateCredit};
    \node[anchor=west, inner sep=0, text width=86mm, align=flush left,
      font=\fontsize{11.5}{17}\selectfont] at (62,43.5) {%
      \quoteText\par\vskip 4mm
      {\color{ruling}\vrule height 0.4pt depth 0pt width 22mm}\par\vskip 2mm
      {\small\itshape\color{muted}\quoteCitation}};
  \end{tikzpicture}}

% === a numbered balloon: a hairline circle with a numeral, centre (#1,#2) in the picture's units ===
\newcommand{\cxballoon}[4]{% x, y, radius in mm, numeral
  \draw[brand, line width=0.7pt, fill=paper] (#1,#2) circle [radius=#3mm];
  \node[text=brand, inner sep=0, font=\fontsize{14}{14}\selectfont] at (#1,#2) {\cxnum{#4}};}

% \treelink{edge style}{label style}{from}{to}{label}: one right angled branch of a tree.
% It leaves the east anchor of the parent, runs level to a junction at two fifths of the gap,
% turns once to the height of the child and arrives level at its west anchor. The first branch
% drawn from a parent carries the shared stretch, so declare the chosen branch first: later
% branches from the same parent begin at the junction, and a dashed branch never overprints the
% heavier path.
\newcommand{\treelink}[5]{%
  \ifcsname cx@trunk@#3\endcsname
    \draw[#1] let \p1=(#3.east), \p2=(#4.west), \p3=($(\p1)!0.4!(\p2)$),
                  \p4=($(\x3,\y2)!0.5!(\p2)$) in
      (\p4) node[#2, above=0.4mm] {#5}
      (\x3,\y1) -- (\x3,\y2) -- (\p2);
  \else
    \expandafter\def\csname cx@trunk@#3\endcsname{}%
    \draw[#1] let \p1=(#3.east), \p2=(#4.west), \p3=($(\p1)!0.4!(\p2)$),
                  \p4=($(\x3,\y2)!0.5!(\p2)$) in
      (\p4) node[#2, above=0.4mm] {#5}
      (\p1) -- (\x3,\y1) -- (\x3,\y2) -- (\p2);
  \fi}

% === list markers: a small square hatched in violet ===
\newcommand{\cxhatchsq}[1]{%
  \raisebox{0.12ex}{\tikz[baseline=0pt]{%
    \fill[hatch fine] (0,0) rectangle (#1,#1);
    \draw[brand, line width=0.55pt, line join=miter] (0,0) rectangle (#1,#1);}}}
\setbeamertemplate{itemize item}{\cxhatchsq{3mm}}
\setbeamertemplate{itemize subitem}{\cxhatchsq{2.1mm}}
\setbeamertemplate{enumerate item}{%
  \tikz[baseline=(enumn.base)]\node[draw=brand, line width=0.7pt, circle, inner sep=0.5pt,
    minimum size=5.4mm, font=\small, text=brand, fill=paper] (enumn) {\cxnum{\insertenumlabel}};}
\setbeamertemplate{enumerate subitem}{\insertenumlabel.}
\setbeamerfont{itemize/enumerate body}{size=\Large}
\setbeamerfont{itemize/enumerate subbody}{size=\large}
\setbeamertemplate{itemize/enumerate body begin}{\vspace*{0mm}}
% beamer resets the list parameters inside \list, so the spacing is set in the list defaults
\makeatletter
\def\@listi{\leftmargin\leftmargini
            \topsep 0\p@ \parsep 0\p@
            \itemsep 6mm \@plus1mm \@minus1mm}
\let\@listI\@listi
\def\@listii{\leftmargin\leftmarginii
             \topsep 0\p@ \parsep 0\p@
             \itemsep 2mm \@plus0.5mm \@minus0.5mm}
\makeatother
% a denser list for reference slides (administration, rosters, small print): a smaller marker and
% the text size #1, \small by default. Set \itemsep right after \begin{cxlist} as with itemize.
\newenvironment{cxlist}[1][\small]{%
  \begingroup
  \setbeamertemplate{itemize item}{\cxhatchsq{2.2mm}}%
  \setbeamerfont{itemize/enumerate body}{size=#1}%
  \begin{itemize}}{\end{itemize}\endgroup}

% === the code plate: mono text in a hairline box under a hatched band ===
% Use it in a [fragile] frame. It is built for lines up to about 60 characters: a longer line
% runs past the box, so shorten the code or the prompt instead of shrinking the type.
\usepackage{listings}
\usepackage[skins,listings]{tcolorbox}
\lstdefinestyle{ainit}{
  basicstyle=\ttfamily\normalsize, keywordstyle=\color{brand}\bfseries,
  commentstyle=\color{muted}\itshape, stringstyle=\color{accent2},
  showstringspaces=false, columns=fullflexible, keepspaces=true, breaklines=false,
  aboveskip=0pt, belowskip=0pt, language=Python}
\newtcblisting{codeplate}[1][]{%
  enhanced, listing only, listing options={style=ainit},
  colback=paper, colframe=ruledark, boxrule=0.6pt, sharp corners,
  left=4mm, right=4mm, top=5.6mm, bottom=2.2mm, before skip=1mm, after skip=1mm,
  overlay unbroken={%
    \fill[hatch] ([yshift=-2.9mm]frame.north west) rectangle (frame.north east);
    \draw[brand, line width=0.5pt]
      ([yshift=-2.9mm]frame.north west) -- ([yshift=-2.9mm]frame.north east);},
  #1}

% === the TikZ styles of the figures ===
\tikzset{
  >={Triangle[length=2.5mm, width=1.9mm]},
  chip/.style={draw=nib, line width=0.7pt, line join=miter, minimum width=1.85cm,
    minimum height=0.82cm, inner sep=3pt, text=ink, fill=paper, font=\large},
  chip strong/.style={chip, draw=brand, line width=1.0pt,
    path picture={%
      \fill[hatch, even odd rule]
        (path picture bounding box.south west) rectangle (path picture bounding box.north east)
        ([shift={(1.25mm,1.25mm)}]path picture bounding box.south west) rectangle
        ([shift={(-1.25mm,-1.25mm)}]path picture bounding box.north east);
      \draw[brand, line width=0.35pt]
        ([shift={(1.25mm,1.25mm)}]path picture bounding box.south west) rectangle
        ([shift={(-1.25mm,-1.25mm)}]path picture bounding box.north east);}},
  chip dim/.style={chip, draw=muted, line width=0.6pt, dash pattern=on 2.4pt off 1.8pt,
    text=muted},
  edge strong/.style={draw=brand, line width=1.3pt, line join=miter,
    -{Triangle[length=2.9mm, width=2.2mm]}},
  edge dim/.style={draw=muted, line width=0.6pt, line join=miter, dash pattern=on 3pt off 2.2pt,
    -{Triangle[length=2.5mm, width=1.8mm]}},
  elab/.style={font=\small\itshape, text=brand, fill=paper, inner sep=1.5pt},
  elab dim/.style={font=\small\itshape, text=muted, fill=paper, inner sep=1.5pt},
  axis/.style={draw=muted, line width=0.6pt, -{Triangle[length=2.3mm, width=1.6mm]}},
  vec a/.style={draw=brand, line width=1.4pt, -{Triangle[length=3.1mm, width=2.3mm]}},
  vec b/.style={draw=accent2, line width=1.4pt, -{Triangle[length=3.1mm, width=2.3mm]}},
  vlab a/.style={text=brand, font=\normalsize},
  vlab b/.style={text=accent2, font=\normalsize},
  anglearc/.style={draw=ink, line width=0.6pt,
    {Triangle[length=2mm, width=1.4mm]}-{Triangle[length=2mm, width=1.4mm]}},
  anglelab/.style={text=ink, font=\normalsize},
  % a note: a numbered balloon on a hairline leader, with its text. The balloons number
  % themselves 1, 2, 3 on each page.
  gloss/.style={font=\small\itshape, text=muted, text width=3.9cm, align=flush left,
    inner xsep=6.9mm, inner ysep=0.6mm,
    execute at begin node={\stepcounter{glossn}},
    path picture={%
      \draw[brand, line width=0.5pt]
        ($(path picture bounding box.north west)+(2.9mm,-5.4mm)$) --
        ($(path picture bounding box.south west)+(2.9mm,0.3mm)$);
      \draw[brand, line width=0.7pt, fill=paper]
        ($(path picture bounding box.north west)+(2.9mm,-2.9mm)$) circle (2.5mm);
      \pgftext[at={\pgfpointadd{\pgfpointanchor{path picture bounding box}{north west}}%
        {\pgfpoint{2.9mm}{-2.9mm}}}]
        {\color{brand}\fontsize{10.5}{10.5}\selectfont\upshape\cxnum{\value{glossn}}}}},
}
      % this file
%   % Copied on 2026-11-03 from eth-ainit-hs26/w1-lecture-material, commit 9792819, file style/ainit-course.tex. To pull a fix, diff against the source and recopy.
% ===== ainit-course.tex: the facts of AINIT HS26 and the three weekends frame =====
% Load this after ainit-preamble.tex in any deck that shows the course roadmap or the title page's
% Ge'ez year. The facts live here once, so a change of date or title is one edit.
%
% Where the facts come from. Dates, themes and titles are the master spreadsheet's Syllabus tab and
% agree with the course website. The three project names differ between the website, the sheet and
% the introduction deck. This file follows the introduction deck (admin-deck/), because that is the
% deck the participants see first. The topics line of each weekend is a short reading of the
% weekend's slots. Re-read the sheet before the semester starts and after any change of schedule.
%
% For the next academic year change \geezYear, the dates, and whatever else moved.

% 2026 in Ge'ez numerals: 20 x 100 + 20 + 6. The title page prints it after the date.
\newcommand{\geezYear}{፳፻፳፮}

% ---- the three weekends (a, b, c): number, title, dates, topics, project, Faust epithet ----
% The epithets come from the course's cold open: Goethe's Faust renders the Greek word logos as
% the Word, the Thought and the Act, and refuses the Power. They map to the three weekends.
\newcommand{\wkaNum}{1}\newcommand{\wkaTitle}{Large language models}
\newcommand{\wkaDates}{16 and 17 October}
\newcommand{\wkaTopics}{Ge'ez games, prompting, OPRO, MLOps}
\newcommand{\wkaProject}{LLM-based marketing agent}\newcommand{\wkaMuse}{the Word}
\newcommand{\wkbNum}{2}\newcommand{\wkbTitle}{Representation learning}
\newcommand{\wkbDates}{13 and 14 November}
\newcommand{\wkbTopics}{Recommendation systems, matrix factorization, bandits}
\newcommand{\wkbProject}{Recommendation systems}\newcommand{\wkbMuse}{the Thought}
\newcommand{\wkcNum}{3}\newcommand{\wkcTitle}{Reinforcement learning}
\newcommand{\wkcDates}{11 and 12 December}
\newcommand{\wkcTopics}{MDPs, Bellman equations, SARSA, REINFORCE}
\newcommand{\wkcProject}{Reinforcement learning for delivery optimization}
\newcommand{\wkcMuse}{the Act}

% ===== frame: the three weekends, three plates side by side =====
% The current weekend has a hatched header band and an ochre flag. Set it with
% \renewcommand{\currentweekend}{b} in a weekend 2 deck. The refused Power hangs as a struck label
% above the gap between the second and the third plate, with a dashed phantom line running down that
% gap. Switch it off with \cxpowerfalse after loading this file.
\newif\ifcxpower
\cxpowertrue
\newcommand{\cxW}{4.08}   % plate width in cm
\newcommand{\cxG}{0.38}   % gap between plates in cm
\newcommand{\cxH}{5.0}    % plate height in cm
\newcommand{\cxtop}{0.56} % height of the row of flags above the plates, in cm
\newcommand{\cxband}{0.44}% height of the header band of a plate, in cm

% the capital form of a Muse: "the Word" becomes "The Word"
\def\cxcaptheA the #1\relax{The #1}
\newcommand{\cxMuse}[1]{{\edef\cxtmp{\csname wk#1Muse\endcsname}\expandafter\cxcaptheA\cxtmp\relax}}

% one plate: weekend letter, left edge in cm, 1 when it is the current weekend.
% The header band carries the weekend label on vellum; on the current plate the rest of the
% band is hatched.
\newcommand{\cxplate}[3]{%
  \begin{scope}[shift={(#2,\cxtop)}]
    \ifnum#3=1
      \fill[hatch] (0,0) rectangle (\cxW,\cxband);
      \def\cxlab{brand}\def\cxout{brand}\def\cxoutw{1.1pt}
    \else
      \def\cxlab{muted}\def\cxout{ruledark}\def\cxoutw{0.65pt}
    \fi
    \node[anchor=west, inner xsep=0.28cm, inner ysep=0, minimum height=\cxband cm, fill=paper,
      text=muted, font=\small\itshape] (lab#1) at (0,{\cxband/2}) {Weekend \csname wk#1Num\endcsname};
    \ifnum#3=1
      \draw[brand, line width=0.5pt] (lab#1.north east) -- (lab#1.south east);
    \fi
    \draw[\cxout, line width=\cxoutw, line join=miter] (0,0) rectangle (\cxW,\cxH);
    \draw[\cxout, line width=0.45pt] (0,\cxband) -- (\cxW,\cxband);
    \ifnum#3=1
      \node[anchor=south west, inner sep=0, fill=accent2, text=paper, font=\small\itshape,
        minimum height=0.46cm, inner xsep=2.6mm] at (0,0) {You are here};
    \fi
    \node[anchor=base west, inner sep=0, text=\cxlab, font=\fontsize{14}{16}\selectfont\itshape]
      at (0.28,{\cxband+0.54}) {\cxMuse{#1}};
    \node[anchor=north west, inner sep=0, text width=3.6cm, align=flush left,
      font=\fontsize{13}{14.4}\selectfont\bfseries] at (0.28,{\cxband+0.82}) {\csname wk#1Title\endcsname};
    \node[anchor=base west, inner sep=0, font=\itshape] at (0.28,{\cxband+2.19})
      {\csname wk#1Dates\endcsname};
    \draw[ruling, line width=0.4pt] (0.28,{\cxband+2.4}) -- (1.1,{\cxband+2.4});
    \node[anchor=north west, inner sep=0, text width=3.6cm, align=flush left,
      font=\fontsize{9}{10.4}\selectfont] at (0.28,{\cxband+2.52})
      {{\color{muted}\csname wk#1Topics\endcsname}\par\vskip 1.2mm
       {\itshape\color{muted}Project:} \csname wk#1Project\endcsname};
  \end{scope}}

\newcommand{\currentweekend}{a}   % the weekend that carries the hatched band and the flag
\newcommand{\threeweekends}{%
  \centering
  \def\cxcura{0}\def\cxcurb{0}\def\cxcurc{0}%
  \expandafter\def\csname cxcur\currentweekend\endcsname{1}%
  \pgfmathsetmacro{\cxxb}{\cxW+\cxG}
  \pgfmathsetmacro{\cxxc}{2*(\cxW+\cxG)}
  \pgfmathsetmacro{\cxxg}{2*\cxW+1.5*\cxG}
  \begin{tikzpicture}[x=1cm, y=-1cm]
    \path[use as bounding box] (0,0) rectangle (13,{\cxtop+\cxH});
    \cxplate{a}{0}{\cxcura}
    \cxplate{b}{\cxxb}{\cxcurb}
    \cxplate{c}{\cxxc}{\cxcurc}
    \ifcxpower
      \node[anchor=base, inner sep=0, text=muted, font=\itshape] (pow)
        at (\cxxg,{\cxtop-0.2}) {The Power};
      \draw[brand, line width=0.8pt] ([shift={(-0.8mm,0.19em)}]pow.base west) --
        ([shift={(0.8mm,0.19em)}]pow.base east);
      \draw[muted, line width=0.5pt, dash pattern=on 2.2pt off 1.8pt]
        (\cxxg,{\cxtop-0.1}) -- (\cxxg,{\cxtop+\cxH-0.1});
    \fi
  \end{tikzpicture}}
        % the facts of the course and the three weekends frame
%   \renewcommand{\deckslug}{AI in Industry, weekend 1}
% and is compiled with lualatex FROM the deck's own directory, three times: the leaf, the dots
% and the balloons are placed with remember picture and settle on the third pass. Fonts are
% EB Garamond, Garamond Math and Libertinus Mono (all in TeX Live) and Kefa III for Ethiopic
% (a macOS system font; set \geezfontpath before this file to use another folder).
%
% What a deck sets:
%   \title \subtitle \author \institute \date       the title page
%   \renewcommand{\deckslug}{...}                   the italic slug in the footer
%   \def\titlegeez{...}                             optional: a Ge'ez line above the cover title
%   \plateImage \plateCredit \quoteText \quoteCitation   the cold open, shown by \coldopenframe
% Frames 'title' and 'cold open' are [plain,noframenumbering]; a \section makes its own divider.
%
% The kit a deck may use:
%   macros   \hl{key words}  \geez{...}  \geezNumeral{n}  \cxballoon  \treelink  \weblink  \qrlink
%            \framedimgh{height}{file}  \coldopenframe  \begin{codeplate}...\end{codeplate}
%            \begin{cxlist}[\small]...\end{cxlist}  a denser list for reference slides
%   styles   chip, chip strong, chip dim, edge strong, edge dim, elab, elab dim, axis, vec a,
%            vec b, vlab a, vlab b, anglearc, anglelab, gloss, hatch, hatch fine
%   colours  paper ink muted brand accent2 good bad ruling ruledark nib
%
% House rules that bind every deck (the lecture-slides skill has the long form):
%   light and projection first; no red against green (good is blue, bad is vermillion); no ALL
%   CAPS; no dash or double hyphen in prose; nothing under a diagram; a section divider carries a
%   numeral, "Section N", the title and a band, nothing else; only the title page names the
%   lecturer; the only number in the chrome is the folio; no cell grids, column letters, tabs,
%   title blocks or arcade chrome. Never name a colour fg or bg: beamer rebinds both.

\usepackage{fontspec}
\usepackage{amsmath}
\usepackage{booktabs}
\usepackage{unicode-math}
\usefonttheme{professionalfonts}
\usefonttheme{serif}
\usepackage{tikz}
\usetikzlibrary{arrows.meta, positioning, calc, backgrounds, fit, overlay-beamer-styles,
                decorations.pathreplacing, patterns.meta}
\setbeamertemplate{navigation symbols}{}

% === type: EB Garamond (TeX Live), Garamond-Math for the maths, Kefa III for Ethiopic ===
\setmainfont{EBGaramond}[
  Extension      = .otf,
  UprightFont    = *-Medium,
  ItalicFont     = *-MediumItalic,
  BoldFont       = *-SemiBold,
  BoldItalicFont = *-SemiBoldItalic,
  Ligatures      = TeX,
  Numbers        = Lining,
  Scale          = 1.08]
\setmathfont{Garamond-Math.otf}[Scale=1.08]
% a monospaced face for code and addresses, sized to the lowercase of the main face
\setmonofont{LibertinusMono-Regular.otf}[Scale=MatchLowercase]

% Kefa III is a macOS system font. Its light and bold instances are variable font instances.
% Without the file, the deck still compiles: Ge'ez text then lacks glyphs and a warning says so.
\providecommand{\geezfontpath}{/System/Library/Fonts/Supplemental/}
\IfFileExists{\geezfontpath KefaIII.ttf}{%
  \newfontfamily\geezfont[Path=\geezfontpath]{KefaIII.ttf}
  % the light instance, for large numerals
  \newfontfamily\geezlightfont[Path=\geezfontpath, RawFeature={instance=Light}]{KefaIII.ttf}
  % the bold instance, for the numerals in balloons (the regular weight is too thin there)
  \newfontfamily\geezboldfont[Path=\geezfontpath, RawFeature={instance=Bold}]{KefaIII.ttf}
}{%
  \PackageWarningNoLine{ainit-preamble}{Kefa III was not found in \geezfontpath. Ge'ez text
    will lack glyphs. Set \string\geezfontpath\space before this file.}%
  \newcommand{\geezfont}{}\newcommand{\geezlightfont}{}\newcommand{\geezboldfont}{}}
\newcommand{\geez}[1]{{\geezfont #1}}
\newcommand{\geezlight}[1]{{\geezlightfont #1}}

% === colours. Never name one fg or bg (beamer rebinds both). ===
\definecolor{paper}{HTML}{F9F4E3}     % vellum, lightness above 93 percent
\definecolor{ink}{HTML}{241F1A}       % warm black, iron gall
\definecolor{muted}{HTML}{5E5448}     % secondary text, darker than black!55
\definecolor{brand}{HTML}{58409A}     % violet in the rubric role
\definecolor{accent2}{HTML}{8F5A06}   % ochre, the second accent
\definecolor{good}{HTML}{0072B2}      % Okabe-Ito blue: a win
\definecolor{bad}{HTML}{D55E00}       % Okabe-Ito vermillion: a loss. No red against green, ever.
\colorlet{ruling}{ink!40!paper}       % the ruled lines of the leaf
\colorlet{ruledark}{ink!62!paper}     % the heavier ruled lines
\colorlet{nib}{ink!88!paper}          % the outline ink of the figures

\setbeamercolor{normal text}{fg=ink, bg=paper}
\setbeamercolor{structure}{fg=brand}
\setbeamercolor{alerted text}{fg=brand}
\setbeamersize{text margin left=15mm, text margin right=15mm}

% === links: violet, and a raw address keeps the mono face, so a link is marked twice ===
% Beamer loads hyperref. The border is cleared; colorlinks stays off because it would also tint
% the section title on the divider.
\hypersetup{pdfborder={0 0 0}, breaklinks=true}
\urlstyle{tt}
\renewcommand{\UrlFont}{\ttfamily\color{brand}}
% \weblink{https://...}{the words the room reads}: an inline link whose visible text is prose.
\newcommand{\weblink}[2]{\href{#1}{\color{brand}#2}}

% \qrlink{2.6cm}{https://...}: a vector QR on a plain white plate in a hairline mount. The code is
% generated from the address at compile time, so it cannot disagree with the address printed
% beside it. Always print the address next to it.
% Two guards, both invisible in the source. nolink: qrcode makes the whole code a link by default,
% and a link is painted in the link colour. \color{black}: qrcode draws in the current colour, so
% on a full bleed frame where \color{white} is active the modules would vanish.
\usepackage{qrcode}
\newcommand{\qrlink}[2]{%
  \tikz{\node[inner sep=4pt, fill=white, draw=ruledark, line width=0.6pt, line join=miter]
    {\color{black}\qrcode[nolink, height=#1]{#2}};}}

% \framedimgh{2.5cm}{img/photo.jpg}: a photograph or figure of the given height in a hairline mount,
% standing 1.4 mm off the picture as the plate of the cold open does.
\newcommand{\framedimgh}[2]{%
  \tikz{\node[inner sep=1.4mm, draw=ruling, line width=0.4pt, line join=miter]
    {\includegraphics[height=#1]{#2}};}}

% === the overlay style used by every piece of page furniture ===
\tikzset{cx page/.style={remember picture, overlay, shift={(current page.north west)},
  x=1mm, y=-1mm}}

% === hatching: the one fill of the figures. Violet lines at 45 degrees, no solid colour under them ===
\tikzset{
  hatch/.style={pattern={Lines[angle=45, distance=0.95mm, line width=0.42pt]}, pattern color=brand},
  hatch fine/.style={pattern={Lines[angle=45, distance=0.62mm, line width=0.36pt]}, pattern color=brand},
}

% === Ge'ez numerals up to five, digits beyond ===
% The room learns ፩ to ፭ in the Ge'ez games (the course numbers the games with them).
\newcommand{\geezNumeral}[1]{\ifcase#1 0\or ፩\or ፪\or ፫\or ፬\or ፭\else\number#1\fi}
\newcommand{\cxnum}[1]{\ifnum#1<6\relax{\geezboldfont\geezNumeral{#1}}\else\number#1\fi}
\newcounter{glossn}

% a box that is scaled down to fit when its natural width exceeds #1, and left alone otherwise
\newsavebox{\cxfitbox}
\newcommand{\cxfit}[2]{%
  \sbox{\cxfitbox}{#2}%
  \ifdim\wd\cxfitbox>#1\relax\resizebox{#1}{!}{\usebox{\cxfitbox}}\else\usebox{\cxfitbox}\fi}

% === the footer slug: each deck renews it ===
\providecommand{\deckslug}{AI in Industry}

% === the leaf: ruled block with a fine double rule, pricking marks in the outer margin ===
\setbeamertemplate{background}{%
  \setcounter{glossn}{0}%
  \begin{tikzpicture}[cx page]
    \draw[ruling, line width=0.55pt] (8,5) rectangle (152,82);
    \draw[ruling, line width=0.25pt] (8.9,5.9) rectangle (151.1,81.1);
    \foreach \i in {0,...,12} {\fill[ruling] (3.2,{5+6.4167*\i}) circle (0.3);}
  \end{tikzpicture}}

% === the title sits above the first ruled line of the block; an over long title is scaled to fit ===
\setbeamertemplate{frametitle}{%
  \vbox to 17.5mm{\offinterlineskip
    \vskip 8.3mm
    \hbox to\textwidth{\vrule height 5.3mm depth 1.9mm width 0pt
      \cxfit{\textwidth}{\color{ink}\fontsize{19}{22}\selectfont\bfseries\insertframetitle}\hfil}%
    \vskip 0.8mm
    \hbox to\textwidth{\color{ruledark}\vrule height 0.7pt depth 0pt width\textwidth}%
    \vskip 0.6mm
    \hbox to\textwidth{\color{ruling}\vrule height 0.3pt depth 0pt width\textwidth}%
    \vss}}

% === the folio in the lower margin, the italic slug beside it ===
\setbeamertemplate{footline}{%
  \vbox to 8.2mm{\offinterlineskip\vss
    \hbox to\paperwidth{\hskip 15mm
      {\color{brand}\fontsize{12}{14}\selectfont\insertframenumber}\hskip 3.4mm
      {\color{muted}\footnotesize\itshape\deckslug}\hfil}%
    \vskip 2.2mm}}

% === the title page ===
% \titlegeez (a Ge'ez line over the title) and \geezYear (the year, in Ge'ez numerals, after the
% date) are optional. \geezYear comes from ainit-course.tex.
\setbeamertemplate{title page}{%
  \begin{tikzpicture}[cx page]
    \node[text=brand, font=\fontsize{17}{20}\selectfont, anchor=base] at (80,20.5)
      {\ifdefined\titlegeez\geez{\titlegeez}\fi};
    \node[text=ink, font=\fontsize{27}{30}\selectfont\bfseries, text width=104mm, align=center,
      anchor=north, inner sep=0] at (80,25) {\inserttitle};
    \draw[ruledark, line width=0.8pt] (62,50.5) -- (98,50.5);
    \draw[ruling, line width=0.3pt] (62,51.5) -- (98,51.5);
    \node[text=muted, font=\large\itshape, anchor=north, inner sep=0] at (80,55)
      {\insertsubtitle};
    \node[text=ink, font=\large, anchor=base, inner sep=0] at (80,67) {\insertauthor};
    \node[text=muted, font=\normalsize\itshape, anchor=base, inner sep=0] at (80,72.2)
      {\insertinstitute};
    \node[text=muted, font=\normalsize\itshape, anchor=base, inner sep=0] at (80,77.2)
      {\insertdate\ifdefined\geezYear\hspace{3.2mm}{\upshape\color{brand}\geez{\geezYear}}\fi};
  \end{tikzpicture}}

% === the section page: a numeral in a heavy square, Section N, the title, a hatched band ===
% The numeral is Ge'ez up to five and a digit beyond. A title that is wider than 92 mm at the
% display size wraps over several lines at a smaller size, and the band moves down under it.
\newif\ifcxwrap
\setbeamertemplate{section page}{%
  \sbox{\cxfitbox}{\fontsize{36}{40}\selectfont\bfseries\insertsectionhead}%
  \ifdim\wd\cxfitbox>92mm\relax\cxwraptrue\else\cxwrapfalse\fi
  \begin{tikzpicture}[cx page]
    \draw[brand, line width=1.8pt, line join=miter] (17,30.2) rectangle (44,57.2);
    \draw[brand, line width=0.4pt] (18.9,32.1) rectangle (42.1,55.3);
    \ifnum\insertsectionnumber<6\relax
      \node[text=brand, font=\fontsize{54}{54}\selectfont, anchor=center] at (30.56,43.7)
        {\geezlight{\geezNumeral{\insertsectionnumber}}};
    \else
      \node[text=brand, font=\fontsize{54}{54}\selectfont, anchor=center] at (30.56,43.7)
        {\number\insertsectionnumber};
    \fi
    \node[text=muted, font=\fontsize{14}{16}\selectfont\itshape, anchor=base west, inner sep=0]
      at (53,33.2) {Section \insertsectionnumber};
    \ifcxwrap
      \node[text=ink, font=\fontsize{26}{29}\selectfont\bfseries, anchor=north west, inner sep=0,
        text width=92mm, align=flush left] (cxst) at (53,36.5) {\insertsectionhead};
      \fill[hatch] ([shift={(0,3.4)}]cxst.south west) rectangle ([shift={(92,7)}]cxst.south west);
      \draw[brand, line width=0.5pt] ([shift={(0,3.4)}]cxst.south west) -- ++(92,0)
        ([shift={(0,7)}]cxst.south west) -- ++(92,0);
    \else
      \node[text=ink, font=\fontsize{36}{40}\selectfont\bfseries, anchor=base west, inner sep=0]
        at (53,47.5) {\insertsectionhead};
      \fill[hatch] (53,53.6) rectangle (145,57.2);
      \draw[brand, line width=0.5pt] (53,53.6) -- (145,53.6) (53,57.2) -- (145,57.2);
    \fi
  \end{tikzpicture}}
\AtBeginSection[]{\begin{frame}[plain,noframenumbering]\usebeamertemplate{section page}\end{frame}}

% === rubrication: key words in semibold violet ===
\newcommand{\hl}[1]{\textcolor{brand}{\bfseries #1}}

% === the cold open: the plate in a hairline mount, the verse, the citation ===
% The deck sets \plateImage (a file name), \plateCredit, \quoteText (lines separated by \\, key
% words in \hl{}) and \quoteCitation, and calls \coldopenframe inside a [plain] frame.
\newcommand{\coldopenframe}{%
  \begin{tikzpicture}[cx page]
    \node[anchor=north west, inner sep=0] (plate) at (16,15)
      {\includegraphics[height=52mm]{\plateImage}};
    \draw[ruling, line width=0.4pt] ([shift={(-1.6,-1.6)}]plate.north west) rectangle
      ([shift={(1.6,1.6)}]plate.south east);
    \node[anchor=north, inner sep=0, text=muted, font=\scriptsize\itshape]
      at ($(plate.south)+(0,3.2)$) {\plateCredit};
    \node[anchor=west, inner sep=0, text width=86mm, align=flush left,
      font=\fontsize{11.5}{17}\selectfont] at (62,43.5) {%
      \quoteText\par\vskip 4mm
      {\color{ruling}\vrule height 0.4pt depth 0pt width 22mm}\par\vskip 2mm
      {\small\itshape\color{muted}\quoteCitation}};
  \end{tikzpicture}}

% === a numbered balloon: a hairline circle with a numeral, centre (#1,#2) in the picture's units ===
\newcommand{\cxballoon}[4]{% x, y, radius in mm, numeral
  \draw[brand, line width=0.7pt, fill=paper] (#1,#2) circle [radius=#3mm];
  \node[text=brand, inner sep=0, font=\fontsize{14}{14}\selectfont] at (#1,#2) {\cxnum{#4}};}

% \treelink{edge style}{label style}{from}{to}{label}: one right angled branch of a tree.
% It leaves the east anchor of the parent, runs level to a junction at two fifths of the gap,
% turns once to the height of the child and arrives level at its west anchor. The first branch
% drawn from a parent carries the shared stretch, so declare the chosen branch first: later
% branches from the same parent begin at the junction, and a dashed branch never overprints the
% heavier path.
\newcommand{\treelink}[5]{%
  \ifcsname cx@trunk@#3\endcsname
    \draw[#1] let \p1=(#3.east), \p2=(#4.west), \p3=($(\p1)!0.4!(\p2)$),
                  \p4=($(\x3,\y2)!0.5!(\p2)$) in
      (\p4) node[#2, above=0.4mm] {#5}
      (\x3,\y1) -- (\x3,\y2) -- (\p2);
  \else
    \expandafter\def\csname cx@trunk@#3\endcsname{}%
    \draw[#1] let \p1=(#3.east), \p2=(#4.west), \p3=($(\p1)!0.4!(\p2)$),
                  \p4=($(\x3,\y2)!0.5!(\p2)$) in
      (\p4) node[#2, above=0.4mm] {#5}
      (\p1) -- (\x3,\y1) -- (\x3,\y2) -- (\p2);
  \fi}

% === list markers: a small square hatched in violet ===
\newcommand{\cxhatchsq}[1]{%
  \raisebox{0.12ex}{\tikz[baseline=0pt]{%
    \fill[hatch fine] (0,0) rectangle (#1,#1);
    \draw[brand, line width=0.55pt, line join=miter] (0,0) rectangle (#1,#1);}}}
\setbeamertemplate{itemize item}{\cxhatchsq{3mm}}
\setbeamertemplate{itemize subitem}{\cxhatchsq{2.1mm}}
\setbeamertemplate{enumerate item}{%
  \tikz[baseline=(enumn.base)]\node[draw=brand, line width=0.7pt, circle, inner sep=0.5pt,
    minimum size=5.4mm, font=\small, text=brand, fill=paper] (enumn) {\cxnum{\insertenumlabel}};}
\setbeamertemplate{enumerate subitem}{\insertenumlabel.}
\setbeamerfont{itemize/enumerate body}{size=\Large}
\setbeamerfont{itemize/enumerate subbody}{size=\large}
\setbeamertemplate{itemize/enumerate body begin}{\vspace*{0mm}}
% beamer resets the list parameters inside \list, so the spacing is set in the list defaults
\makeatletter
\def\@listi{\leftmargin\leftmargini
            \topsep 0\p@ \parsep 0\p@
            \itemsep 6mm \@plus1mm \@minus1mm}
\let\@listI\@listi
\def\@listii{\leftmargin\leftmarginii
             \topsep 0\p@ \parsep 0\p@
             \itemsep 2mm \@plus0.5mm \@minus0.5mm}
\makeatother
% a denser list for reference slides (administration, rosters, small print): a smaller marker and
% the text size #1, \small by default. Set \itemsep right after \begin{cxlist} as with itemize.
\newenvironment{cxlist}[1][\small]{%
  \begingroup
  \setbeamertemplate{itemize item}{\cxhatchsq{2.2mm}}%
  \setbeamerfont{itemize/enumerate body}{size=#1}%
  \begin{itemize}}{\end{itemize}\endgroup}

% === the code plate: mono text in a hairline box under a hatched band ===
% Use it in a [fragile] frame. It is built for lines up to about 60 characters: a longer line
% runs past the box, so shorten the code or the prompt instead of shrinking the type.
\usepackage{listings}
\usepackage[skins,listings]{tcolorbox}
\lstdefinestyle{ainit}{
  basicstyle=\ttfamily\normalsize, keywordstyle=\color{brand}\bfseries,
  commentstyle=\color{muted}\itshape, stringstyle=\color{accent2},
  showstringspaces=false, columns=fullflexible, keepspaces=true, breaklines=false,
  aboveskip=0pt, belowskip=0pt, language=Python}
\newtcblisting{codeplate}[1][]{%
  enhanced, listing only, listing options={style=ainit},
  colback=paper, colframe=ruledark, boxrule=0.6pt, sharp corners,
  left=4mm, right=4mm, top=5.6mm, bottom=2.2mm, before skip=1mm, after skip=1mm,
  overlay unbroken={%
    \fill[hatch] ([yshift=-2.9mm]frame.north west) rectangle (frame.north east);
    \draw[brand, line width=0.5pt]
      ([yshift=-2.9mm]frame.north west) -- ([yshift=-2.9mm]frame.north east);},
  #1}

% === the TikZ styles of the figures ===
\tikzset{
  >={Triangle[length=2.5mm, width=1.9mm]},
  chip/.style={draw=nib, line width=0.7pt, line join=miter, minimum width=1.85cm,
    minimum height=0.82cm, inner sep=3pt, text=ink, fill=paper, font=\large},
  chip strong/.style={chip, draw=brand, line width=1.0pt,
    path picture={%
      \fill[hatch, even odd rule]
        (path picture bounding box.south west) rectangle (path picture bounding box.north east)
        ([shift={(1.25mm,1.25mm)}]path picture bounding box.south west) rectangle
        ([shift={(-1.25mm,-1.25mm)}]path picture bounding box.north east);
      \draw[brand, line width=0.35pt]
        ([shift={(1.25mm,1.25mm)}]path picture bounding box.south west) rectangle
        ([shift={(-1.25mm,-1.25mm)}]path picture bounding box.north east);}},
  chip dim/.style={chip, draw=muted, line width=0.6pt, dash pattern=on 2.4pt off 1.8pt,
    text=muted},
  edge strong/.style={draw=brand, line width=1.3pt, line join=miter,
    -{Triangle[length=2.9mm, width=2.2mm]}},
  edge dim/.style={draw=muted, line width=0.6pt, line join=miter, dash pattern=on 3pt off 2.2pt,
    -{Triangle[length=2.5mm, width=1.8mm]}},
  elab/.style={font=\small\itshape, text=brand, fill=paper, inner sep=1.5pt},
  elab dim/.style={font=\small\itshape, text=muted, fill=paper, inner sep=1.5pt},
  axis/.style={draw=muted, line width=0.6pt, -{Triangle[length=2.3mm, width=1.6mm]}},
  vec a/.style={draw=brand, line width=1.4pt, -{Triangle[length=3.1mm, width=2.3mm]}},
  vec b/.style={draw=accent2, line width=1.4pt, -{Triangle[length=3.1mm, width=2.3mm]}},
  vlab a/.style={text=brand, font=\normalsize},
  vlab b/.style={text=accent2, font=\normalsize},
  anglearc/.style={draw=ink, line width=0.6pt,
    {Triangle[length=2mm, width=1.4mm]}-{Triangle[length=2mm, width=1.4mm]}},
  anglelab/.style={text=ink, font=\normalsize},
  % a note: a numbered balloon on a hairline leader, with its text. The balloons number
  % themselves 1, 2, 3 on each page.
  gloss/.style={font=\small\itshape, text=muted, text width=3.9cm, align=flush left,
    inner xsep=6.9mm, inner ysep=0.6mm,
    execute at begin node={\stepcounter{glossn}},
    path picture={%
      \draw[brand, line width=0.5pt]
        ($(path picture bounding box.north west)+(2.9mm,-5.4mm)$) --
        ($(path picture bounding box.south west)+(2.9mm,0.3mm)$);
      \draw[brand, line width=0.7pt, fill=paper]
        ($(path picture bounding box.north west)+(2.9mm,-2.9mm)$) circle (2.5mm);
      \pgftext[at={\pgfpointadd{\pgfpointanchor{path picture bounding box}{north west}}%
        {\pgfpoint{2.9mm}{-2.9mm}}}]
        {\color{brand}\fontsize{10.5}{10.5}\selectfont\upshape\cxnum{\value{glossn}}}}},
}
